\begin{center}
{\Large\bfseries Hecke points, theta lifts and residual irreducibility}\\[4pt]
Edition 0.86 \quad 25 November 2026
\end{center}

Fix one of the original members in the denominator family. Its upper
Hecke labels have rank $r+t$ and its lower theta rows have rank $r<t$.
We construct the Galois representations of the upper points, identify
the single upper label containing a nonzero lifted row, and recover
the lower parameter by removing a character string of length $t$.
The argument at the two type primes controls the residual eigenvalues
on the specified inertia elements. It proves irreducibility and descends
the lower representation to the coefficient ring prescribed by the
member. The final comparisons record which of these statements hold on
the original family and which multiplicity and selection questions remain.

\begingroup\emergencystretch=3em
\section{The first ten conditions at a coefficient row}

Fix an original member $D=w.D(Z)$ and one of its coefficient indices $i$.
Write $E/F$ for its CM extension, $c$ for complex conjugation,
$\ell=\ell_i$ and $\iota=\iota_i$ for the chosen complex embedding.
The member carries both a realization of its cuspidal labels and maps reading
their local parameters. B85U1 identifies the cuspidal labels with their
actual automorphic representations, injectively on classes of a fixed rank.
M84T1 identifies the period readings and M84T6 the character frame.
We use these three parts of the member in what follows.

% BEGIN S86A1
\paragraph{S86A1. Recovering the global list from its local factors.}
For a triple $(a,\pi,d)$ of positive ranks, expand its local
Speh parameter into the $d$ translates of $\pi$ by
\[
              (d-1)/2,\ (d-3)/2,\ \ldots,\ (1-d)/2.
\]
If two lists give equal parameters away from finitely many primes, the
standard unramified polynomials of the expanded lists agree there.
By Jacquet--Shalika uniqueness in A72F5 the two expanded multisets of
cuspidal classes with real norm twists agree. Each base cusp is unitary,
so its real twist exponent is recovered from the absolute value of the
central character (the norm exponent on the center is multiplied by
$a$). Removing the twist then determines its base cusp and hence its
original label. For each base cusp separately, let $m(s)$ count
occurrences at the nonnegative half-integer $s$. The number of strings
of length $2s+1$ is exactly $m(s)-m(s+1)$: all longer strings of
the same parity occur at both exponents, and the other parity occurs
at neither. This recovers ranks, labels and lengths with all repetitions.
For $d=1$ it gives ordinary cuspidal multiplicity one.
These are the two list comparisons in the selected row.

Split Frobenius elements away from any finite excluded set are dense in
$G_E$ when all lifts of places are allowed. Given a coset in a finite
Galois quotient over $E$, take a common normal closure over $F$. Lift
the requested element to the subgroup fixing $E$. Chebotarev over $F$
gives its conjugacy class outside the excluded primes, and choosing a
place above it gives this very lift. Its restriction to $E$ is trivial,
so the prime splits in $E/F$ and the residue norm at its lift in $E$
is unchanged. This meets the given coset. For the last clause of
$\mathrm{adm5}$ use the Brauer--Nesbitt theorem for semisimple
finite-dimensional representations with the same characteristic
polynomials on group elements, over the specified coefficient field.
In a basis its module isomorphism is the required simultaneous
conjugacy. This proves $\mathrm{an1},\mathrm{an2},\mathrm{a5}$.
% END S86A1

% BEGIN S86A2
\paragraph{S86A2. Characters, Satake roots and the local labels.}
For two rank-one labels with equal readings at almost all split places,
take the quotient $\xi$ of the unitary Hecke characters read by the member.
At split unramified places the inverse uniformizer values are precisely
the single roots of the given parameters. Hence $\xi$ is trivial there.
The finite-place Euler products for $L_E(s,\xi)$ and $\zeta_E(s)$ have a
quotient holomorphic and nonzero near $1$. Indeed at an inert prime
over $v$ their local logarithms have difference bounded by
$O((\mathrm Nv)^{-2\Re s})$ on compact sets in $\Re s>1/2$.
At the remaining finite exceptional primes the factors are regular and
nonzero near $1$ by unitarity. Consequently $L_E(s,\xi)$ has a pole
at $1$. The Hecke pole theorem (Tate's thesis, global zeta integral)
implies $\xi=1$, also for unitary characters with a norm component.
Injectivity of the character reading proves $\mathrm{HeckeSplitOne}$.

For each unramified positive-rank cuspidal label the Jacquet--Shalika
bound of A72F5 and its contragredient give
$q^{-1/2}<|\alpha|<q^{1/2}$ for all its Satake roots. The source
reads inverse roots, which satisfy the same pair of inequalities.
At all places a unitary character has bounded Weil image and zero
monodromy. This proves both clauses of $\mathrm{temp}$, even at
ramified places in rank one. The characters $\eta,\eta^{-1}$ are
conjugate self-dual: the defining restriction of $\eta$ to the real
subfield is its quadratic class character, trivial on norms, so
$\eta(x)\eta(c x)=1$. Their infinity multisets are the opposite
singletons $\{\pm1/2\}$ in the two embeddings, by the character
frame, and they and every $\mu_i$ are unitary. This gives
$\mathrm{TieAut}$.

At a place of a type prime $g$ the member reads the monic polynomial
whose roots lifted by its Teichm\"uller character are the eigenvalues of
the designated tame generator on the wild-invariant part of the
Weil space. A real norm twist is trivial on inertia, and this subspace
has dimension at most the full rank. Thus the label is unchanged by
the twist and its degree is bounded by the rank. This gives
$\mathrm{tieLab}$. Complete splitting of the type primes makes each
chosen place split over $F$, giving $\mathrm{gsp}$.
% END S86A2

% BEGIN S86A3
\paragraph{S86A3. Sums, twists and the original period readings.}
In the Weil--Deligne model the operations of direct sum, unramified
twist and tensor with a character are literal operations on
representations. Tensor distributes over direct sum, composes by
character multiplication and sends an inertia eigenvalue $z$ to
$\chi(\gamma)z$. The inertia polynomial is monic of degree the rank.
The trivial unramified line acts identically. Transporting through the
member's identification proves every identity in $\mathrm{TieLP}$.

At $w\mid\ell$, a de Rham representation becomes semistable over
a finite extension by the $p$-adic monodromy theorem (A72F2).
On its semistable period space $\ker N$ is the crystalline period
space there. Thus it is potentially crystalline exactly when it is
de Rham and $N=0$. It is crystalline over the original field exactly
when, in addition, the finite inertia descent action is trivial:
first descend the unramified constants by semilinear Frobenius
descent, while the ramified descent acts through inertia on the
period basis. Equivalently the resulting Weil--Deligne parameter is
unramified. Conversely a crystalline basis has these properties.
In a de Rham comparison each labelled filtered vector space has
dimension the rank, so the dimensions of its successive graded
quotients add to that rank, with repeated jumps counted separately.
The member's de Rham, jump and Weil--Deligne readings are those of
this comparison by its period frame and M84T1.

Away from $\ell$, write the inertia action in the form
$r(\gamma)\exp(t_\ell(\gamma)N)$ with $r$ of finite inertia
image. Frobenius semisimplification preserves this inertia and $N$.
An isomorphism of the parameter with a direct sum therefore gives
complementary inertia-stable spaces of precisely the two specified
dimensions. On a summand with $N=0$ a finite-order inertia matrix
in characteristic zero is semisimple. A characteristic polynomial
$(X-c)^m$ makes it scalar $c$ there, including after transport by
$\iota^{-1}$. Conjugation by a basis change transports each period
isomorphically. Direct sum takes direct sums of invariants and
filtrations. Tensor with a de Rham line shifts the labelled
filtration by that line's sole jump. Its inverse gives the converse
de Rham implication. These operations give the same sums and twists
of the associated Weil--Deligne parameters, away from $\ell$ by the
inertia formula and above $\ell$ by the semistable comparison after
one common finite extension. This supplies $\mathrm{TieGal}$ with its
de Rham guards. Together with S86A2 it proves $\mathrm{tie0}$.
% END S86A3

% BEGIN S86A4
\paragraph{S86A4. The cuspidal Galois attachment on this member.}
Fix a regular algebraic conjugate-self-dual cuspidal label $\pi$ of
rank $a\ge2$. Its realization from B85U1 has exactly these
properties, and its polarization character can be taken to be the
$a$th power of the quadratic character of $E/F$. This character is
trivial on norms and has the required sign $(-1)^a$ at infinity.
The contragredient identification gives the polarization of A72F5--8.
Take its single continuous semisimple $\ell$-adic attachment. It is
de Rham at \emph{every} $w\mid\ell$, and its jumps for an embedding
over $w$ are $(a-1)/2-z$, with $z$ running through the labelled
infinity multiset. We retain this unconditional de Rham conclusion
of the attachment theorem for use below.

The finite normalization is $\pi_w|\det|^{(1-a)/2}$. Both
Weil spaces have dimension $a$. At spherical places away from
$\ell$ the parameters agree. Under the source's additional inertia
test $\mathrm{Nc}$ the automorphic parameter is tame, has zero
monodromy and satisfies the stated repeated-root condition on the
chosen tame generators. An unramified twist preserves all three
properties. The member's cuspidal realization has good supercuspidal support
at every finite place (\(\mathrm{CGTrueT.good}\) in B85U1). Good
support is preserved by a norm twist. The good-support comparison
in A72F7 thus identifies the tested correspondence with the
classical one on these inputs over every completion. Apply the
polarized attachment and local comparison of A72F8. This gives
zero monodromy and equal inertia polynomials in the guarded case
away from $\ell$, and the zero-monodromy and inertia assertions of
$\mathrm{CuspA}$ in the guarded case above $\ell$. In the latter
case the domination of monodromy by a zero-monodromy parameter
forces zero monodromy on the Galois side as well. The equalities
are on the member's actual $\mathrm{wdOf}$ via its period
identification. When $a=1$, global reciprocity and the local character
comparison of A72F3 give the equality of the entire parameter at
every finite place and de Rhamness at every place above $\ell$.

Two further conclusions belong to this same field. For every regular
algebraic cuspidal label the conjugate infinity multiset is
$\{w_0-z\}$ by purity in A72F5. For every rank-one label $\chi$
and every real $s$ making all $z+s$ integral, the character
$\chi|\cdot|^s$ is algebraic. Its $\ell$-adic avatar in the source's
inverse-uniformizer convention has the indicated finite-part
parameter everywhere and labelled jumps $-(z+s)$ above $\ell$
by A72F1--3. This proves the entire $\mathrm{a2}$ field, including
these characters which were not assumed regular before the shift.
Finally the field called $\mathrm{g4p}$ in this version quantifies
the proposition $\mathrm{True}$ at each $i$ and is satisfied at once.
Together S86A1--4 establish all ten entries of the first part
$\mathrm{RowsAAlt}$ of $\mathrm{what}_2$ on the original member.
% END S86A4

\section{Galois representations of the full Hecke points}

% BEGIN H86P1
\paragraph{H86P1. An actual base-change list at every localized label.}
The host satisfies $E=FK$ for $K=\mathbb Q(\sqrt{-d})$: write
$x=(x+cx)/2+\sqrt{-d}\,(x-cx)/(2\sqrt{-d})$ with both coefficients
fixed by $c$. Its prescribed real degree is at least two, and its CM
extension is unramified at finite places by the host data. Every
place above $2$ splits. Indeed $d+1=8m$ and the conjugate integers
$(1\pm\sqrt{-d})/2$ have sum $1$ and product $2m$.
Exactly one is in a given prime above $2$, which is therefore moved
by conjugation. The member's positive definite Hermitian spaces give
definiteness for all rows $1\le j\le m_{\rm row}$ in the design range.

Now fix $P\in\mathrm{AV}_i$, including any representative with
no Hecke point. The member's base-change datum embeds it in the
unitary labels of row $j(i)$ and gives sphericity there at every
nonsplit place. The positive rank is $n=r+t$. Apply B85U2 to this
image with the field hypotheses just checked. The embedding preserves
the split local parameter, the nonsplit spherical tuple and the
infinity multiset. Consequently the \emph{same} list $\mathcal P$
satisfies all six clauses of $\mathrm{IsBCV}(i,P,\mathcal P)$.
This proves $\mathrm{adm1}$ on the selected member.
% END H86P1

% BEGIN H86P2
\paragraph{H86P2. The representation attached to a list and its weights at every coefficient place.}
Fix any list from H86P1, including an alternative list satisfying
the same tests. For each $(a,\pi,d)$ in it and $0\le k<d$, put
$e=d-2k-n$. The $k$th term of the Speh parameter after overall
normalization is $\pi|\det|^{e/2}$. If $z$ is an exponent of
$\pi$ at a complex embedding, then $z+(d-1)/2-k$ occurs in the
given infinity multiset of $P$. This multiset is regular and lies
in $(n-1)/2+\mathbb Z$ by its dominant coefficient weight.
Thus $z+e/2\in\mathbb Z$ and the exponents within the term
are distinct.

Here is the half-integral twist in the Galois construction. Write
$u=(e+a-1)/2$. If $u\in\mathbb Z$, $\pi$ is regular algebraic
in rank $a$. Take the usual attachment, normalized by $(1-a)/2$,
and tensor with the algebraic character for $|\cdot|^u$ in the
convention of S86A4. If $u\in\mathbb Z+1/2$, first replace
$\pi$ by $\pi\eta$. At each embedding the exponent of $\eta$
is $1/2$ or $-1/2$. Thus the new exponents lie in
$(a-1)/2+\mathbb Z$, still without repetition. This twist preserves
cuspidality and conjugate self-duality. Tensor its attachment with
the avatar of $\eta^{-1}|\cdot|^u$, an algebraic character since
its infinity exponents are integral. In both cases the auxiliary
character has unramified smooth part at every finite place. The
normalized finite parameter has exactly the requested twist, and
the jumps are $-(z+e/2)$. For example in the second case the sum of
the jumps of the factors is
$(a-1)/2-(z+h)-(-h+u)=-z-e/2$, with $h$ the exponent of $\eta$.

Taking the direct sum over all terms gives a continuous semisimple
representation $R_{\mathcal P}$ of dimension $n$. At each place
where the list is spherical, every term has unramified parameter
and zero monodromy: inertia and the nilpotent operator are
block-diagonal on the sum. Away from $\ell$ its parameter is the
normalized list at these places, and above $\ell$ it is crystalline
there by the spherical case of A72F8 or A72F3. At every $w\mid\ell$,
including the two type primes when applicable, every factor is de
Rham by the full attachment theorem and the rank-one comparison.
For $\tau$ inducing $w$, the jumps are
\[
        \bigl\{(n-1)/2-z':z'\in\mathrm{infWt}(i,\iota\tau)\bigr\}
\]
with multiplicity. In particular these jumps are distinct at every
embedding. This application uses the full de Rham assertion for
regular algebraic cusps even at places where the source asks only
for a conditional Weil--Deligne comparison.
% END H86P2

% BEGIN H86P3
\paragraph{H86P3. Existence and uniqueness at a point.}
For $x\in X_i$ take any $P$ above $x$ by the member's point
surjection and any base-change list by H86P1. At a split place
outside $S_i,Q_i$ and $\ell$, the original level gives a spherical
parameter. Reading the Frobenius of H86P2 gives the original Hecke
polynomial of $x$, by the split-point identity including its norm
normalization. Any lift of that Frobenius has the same reading by
unramifiedness. Thus $R_{\mathcal P}$ is a rational point
representation. Two semisimple point representations have equal
characteristic polynomials on this dense set by S86A1. Continuity
of the coefficients gives equality everywhere, and Brauer--Nesbitt
gives an isomorphism. This includes the scalar extension of every
integral good point representation: semisimplicity persists over
an extension of a characteristic-zero field, by separability of the
semisimple image algebra. Conversely the member's original large
coefficient duty (M84T6) realizes each rational point representation
up to conjugacy over $O_i$. Its matrices are continuous for the
valuation topology of this ring, the original Frobenius identities
descend by injectivity and semisimplicity descends to its fraction
field. Choosing at each point gives the required integral family.

Two labels above the same point consequently give isomorphic list
representations. At all but finitely many places away from $\ell$,
both lists are spherical, and H86P2 and cancellation of the common
unramified twist identify their parameters. S86A1 recovers the
original lists up to permutation, including the Speh lengths.
Their parameters hence agree at \emph{all} places. The two
$\mathrm{IsBCV}$ readings give the requested equal split parameters
and equal base changes of the nonsplit Satake tuples.
% END H86P3

% BEGIN H86P4
\paragraph{H86P4. The four integral family properties.}
Fix any integral good family $R_x$ at $i$. For each $j$ and each
$g\in G_E$, the $j$th coefficient of the reversed characteristic
polynomial of $R_x(g)$ is a function of $x$. This function belongs to
the original integral Hecke algebra $T_i\subset O_i^{X_i}$.
Indeed it does so on the split good Frobenius elements by the
defining Hecke polynomial, and is continuous in $g$ in the finite
product topology. The algebra is closed there: as an $O_i$-submodule
of a finite free module it is finitely generated, and is the image
of a compact finite power of $O_i$. Use S86A1's density. This is
$G1$. Residual evaluations at different points of the same $T_i$
agree. Each evaluation followed by reduction is surjective by the
constants and has maximal kernel, necessarily the unique maximal
ideal of $T_i$. Since each map fixes the residue-field constants
their induced maps coincide. This proves independence of the full
residual characteristic polynomial at every $g$.

At $w\nmid\ell$ outside $S_i\cup Q_i$ the local list of each
point is spherical, by the split level or the inert base change.
H86P2--3 give unramifiedness in the rational model. Injectivity
of $O_i$ into the coefficient field gives trivial inertia already
integrally, at every conjugate choice of place. This is $G4$.
At $w\mid\ell$, when $\ell$ is neither type prime, such a list is
also spherical: $Q_i$ excludes these places, and the chosen
auxiliary prime is distinct from $\ell$ even when $\ell$ is the
first auxiliary prime. The split level and inert reading exhaust
the possibilities. H86P2--3 give crystallinity and the stated
labelled jumps, proving $G2a$.

Finally take one finite splitting extension of $\operatorname{Frac}(O_i)$
for the representations at every $x\in X_i$. For instance, split the
semisimple envelope of their direct sum over an algebraic closure and
adjoin the finitely many coefficients of a splitting basis and its
matrix units. All simple submodules are now absolutely irreducible.
If two disjoint simple submodules of $R_x$ were isomorphic, after
scalar extension they would give a positive-dimensional repeated
direct summand. At any one embedding above $\ell$, de Rham periods
of a direct sum are the sum of the periods. A direct summand of a
de Rham representation is de Rham, and its labelled jump multiset
has positive cardinality equal to its dimension. Every one of its
jumps would occur at least twice in the ambient multiset, contrary
to H86P2. This proves the distinctness assertion also when $\ell$ is
a type prime and the source imposes its additional activity guard.
Together with residual polynomial independence this is $G3$.
H86P3 and these four assertions prove the full
$\mathrm{provisionFortyOne}$ for every original $i$.
Of the three conjuncts of the selected field $\mathrm{Conclusion1}$,
the middle one is the trivial $\mathrm{g4p}$ of S86A4. Only
$\mathrm{Conclusion42}$, the parameter assertion for nonzero
theta rows, remains for that field.
% END H86P4

\section{Down-up scalars and ranks on the original member}

% BEGIN L86M1
\paragraph{L86M1. The level restriction of the theta operator.}
Fix a member $D=w.D(Z)$ admitted by the original layer point. Its
realization supplies the pinned characterisation, primitive conditions,
and the two readings $Y,Z_L$. Take $1\le j\le m$, write $j=i+1$ and
fix a used tuple $\mu$, either $\mu_0$ or $\mu_{\rm At}(w')$.
Choose $t=Y.\mathrm{tupT}(\mu)$, so $Y.\mathrm{tup}(t)=\mu$.
Write $V,W$ for the source and target Hilbert models at $i,t$ and
the distinguished embedding, $F=R.\mathrm{thG}(i,t):V\to W$, and
$\theta:M_t(\mu,j)\to M_j$ for the member's down lift.
Use the native coefficient embedding on both sides. If the place
coordinate uses its conjugate, transport both spaces, actions and
characters by the same conjugation. The pin gives $R.\mathrm{phs}(i)=1$ in the conjugated formula for the lift.
Inner-product identities and real eigenvalues survive this transport.

Put $P=P_0$, $N=|P|$, and $\xi(h)=\mathrm{emb}_{LL}(\psi_0(h)^{-1})$.
The subgroup acts unitarily on $V,W$. On $V$ the action extends the
action on rational forms, which preserves the class-sum pairing.
Those forms span $V$ by the Hilbert-model basis. On $W$ unitarity is
the member's primitive group condition. The primitive equivariance
identity for $F$ holds on the same rational spanning set, hence on
$V$. On either space put
\[
       E=\frac{1}{N}\sum_{h\in P}\xi(h)^{-1}\rho(h).
\]
Translation of the summation variable identifies its image with the
$\xi$-space. Unitarity, $|\xi(h)|=1$ and inversion in $P$ give
$E^*=E$. In particular $FE_V=E_WF$ and $E_VF^*=F^*E_W$.

There are maps $A:M_t(\mu,j)\to V$ and $B:M_j\to W$ with
\[
 \operatorname{im} A=E_V V,\qquad
 \operatorname{im} B=W_0\subset E_W W,\qquad
 A^*A=N\,1,\quad B^*B=N\,1,\qquad B\theta=FA.
\]
Indeed, the native source has the basis obtained from the rational
level columns by $j_{M_t}e_{M_t}$. Map them by the level inclusion
followed by $j_V$. The inclusion columns form a basis for the image
of the subgroup average. If $C$ is their matrix and $H$ the uncut
Hermitian form, the source metric is $N^{-1} C^*HC$. This proves
the image and metric assertions for $A$. On the target take
$B=Y.\mathrm{mj}.eW_i$. The target reading gives its squared-norm
factor $N$ and its image, the base-level space $W_0$; the
primitive level condition puts this space in the $\xi$-space.
Here $W_0$ also equals the span of the base-level target columns:
the auxiliary level equality identifies the model space with that
span, and scalar extension of the column basis gives the same range.
For the last identity it suffices to check a rational source column.
Its native and model readings evaluate the same theta integral
on pure tensors in the lowest compact types. The source realization
is injective, including when $t$ is replaced by its chosen tuple
representative. The difference of the two equivariant target
multiplicity maps vanishes on those tensors. Each lowest type
generates its irreducible discrete-series factor, so the difference
vanishes on the generated tensors and, by continuity, on the full
smooth model. The target evaluation is injective, which proves the
claimed equality on the column and hence everywhere.
% END L86M1

% BEGIN L86M2
\paragraph{L86M2. The restricted adjoint and the dimension of every lift image.}
For any $y\in M_j$ we have $E_WBy=By$. L86M1 gives
$F^*By\in E_VV=\operatorname{im} A$. For all $x\in M_t(\mu,j)$,
\[
 \langle Ax,F^*By\rangle
 =\langle B\theta x,By\rangle=N\langle\theta x,y\rangle
 =\langle Ax,A\theta^*y\rangle .
\]
Both vectors tested against $\operatorname{im}A$ lie in that space,
so $A\theta^*y=F^*By$. Therefore
$B\theta\theta^*y=FF^*By$. The native up lift is $\theta^*$.
For every parameter part $W_{ip}$, the primitive scalar identity
says $FF^*z=\lambda_{itp}z$ on $W_{ip}$, with
$\lambda_{itp}=R.\mathrm{lam3}(i,t,p,v_0)>0$. The primitive
target decomposition is
\[
 W_0=\bigoplus_p (W_{ip}\cap W_0).
\]
On the inverse image under $B$ of each summand,
$\theta\theta^*$ is multiplication by $\lambda_{itp}$, since $B$
is injective. The inverse on this finite sum multiplies its $p$-th
component by $\lambda_{itp}^{-1}$. It follows that $\theta^*$ is
injective and that its image has dimension $\dim M_j$, independently
of $\mu$. For a twisted place the tuple identity
$\mu_{\rm At}(w')\in\{\mu_0,\mu_{\rm Tw}(w')\}$ and the twisted-place
guard give $\mu_{\rm Tw}(w')=\mu_{\rm At}(w')$. Applying the image
dimension formula to it and to $\mu_0$ proves the original member's
rank comparison at every allowed $j$, for every native size datum.
% END L86M2

% BEGIN L86M3
\paragraph{L86M3. One family part for a finite type.}
Let $\sigma$ be a native finite type. Its reading is an isomorphism
class of irreducible smooth finite-adelic representations. Choose a
representative $\pi$; the conjugate row representation occurs at one
of the two row levels. At a good spherical place $v$ put $K$ equal
to its hyperspecial group. Each spherical double-coset operator
$T_d$ has a single scalar $c(v,d)$ on $\pi^K$. To see this, take a
nonzero vector in $\pi^K$. Its full orbit spans $\pi$ by
irreducibility. Averaging the orbit over $K$ shows that its $K$-fixed
space is generated by the spherical double-coset sums together with
translations at the complementary places. This holds for every nonzero
fixed vector. These operators thus act irreducibly on $\pi^K$ and
$T_d$ commutes with them, by commutativity of the unramified spherical
algebra and of the distinct local factors. The smooth cyclic
representation has countable dimension: an open stabilizer has
countably many translates in the second-countable finite-adelic
group. On a countable-dimensional simple complex module a commuting
operator $S$ is scalar. Indeed, if every $S-a$ is invertible then,
for fixed $u\ne0$, the uncountably many $(S-a)^{-1}u$ have a
nontrivial finite dependence. Clearing distinct denominators yields
a nonzero polynomial in $S$ killing $u$, although each of its linear
factors is invertible. Otherwise simplicity gives $S=a$ for some
$a$. Where the fixed space vanishes the scalar can be chosen
arbitrarily. At unused indices extend $c$ arbitrarily.

The good-level vectors of an isotypic space lie in the joint
eigenspace: each is a finite sum from equivariant images of $\pi$;
averaging its preimages over $K$ fixes the sum and makes each term
$K$-fixed. The pinned standard frames transfer the same
double-coset sums to the lift and row models. The row model uses
$\bar\pi$, giving eigenvalue $\overline{c(v,d)}$. Inversion gives
the adjoint double coset $d^\vee$ for the invariant inner product:
the two finite coset sums are convolution by the characteristic
functions of inverse double cosets, with the same normalisation.
On any nonzero occurrence this gives
$\overline{c(v,d)}=c(v,d^\vee)$, the convention in the defining
row eigenspace $E(c)$. Thus some nonzero row vector lies in
$E(c)$ at one of the two levels, and the lift reading of every
$\sigma$-vector at $j=i+1$ belongs to $W_i(c)$.

Write $p$ \emph{carries} $c$ if $E(c)\ne0$ and the adjoint of the
reference theta map takes $E(c)$ into the span of rational vectors of
inert type $p$. This condition depends only on $E(c)$. The $p$-th
row and lift parts are respectively the sums of $E(c')$ and
$W_i(c')$ for the systems $c'$ carried by $p$. If no $p$ carried
$c$, every nonzero such $E(c')$ would differ from $E(c)$ on some
tested double-coset eigenvalue. Otherwise the defining systems
of equations give the same subspace, which has the same carriers.
Distinct eigenvalues of one normal spherical operator give orthogonal
vectors; normality follows because its inverse-coset adjoint is in
the same commutative algebra. Our nonzero row vector would be
orthogonal to every row part. But at its own level the row space is
the sum of the intersections with these parts. This contradicts
positive definiteness. Choose a carrier $p=p(\sigma)$; for every
$i$ we have $B(\mathrm{isot}_j(\sigma))\subset W_{ip}$.
% END L86M3

% BEGIN L86M4
\paragraph{L86M4. The original field \textnormal{\texttt{lam}}.}
On a nonzero native isotypic summand L86M2 and L86M3 give
$\theta\theta^*y=\lambda_{it,p(\sigma)} y$. The member's reading
for $\mathrm{lamP}$ identifies its named real scalar with any real
scalar acting by this formula on that nonzero summand. It follows
that $\mathrm{lamP}(\mu,j,\sigma)=\lambda_{it,p(\sigma)}>0$.
On a zero summand the operator equation holds at its only vector
and the positivity implication has false antecedent. This covers
every used tuple and every allowed index. In particular a location
has $\mu=\mu_{\rm Tw}(w')=\mu_{\rm At}(w')$ by its twisted-place
guard, so the same calculation gives the location-wise field.
% END L86M4

% BEGIN L86M5
\paragraph{L86M5. Independence and the remaining dimension comparison.}
In the common ambient multiplicity representation the full
$\pi$-isotypic space is the algebraic sum of images of equivariant
maps out of $\pi$. Each nonzero image is simple. For any finite sum
of simple images one may discard an image contained in the preceding
sum; every other image meets that preceding sum trivially. The
resulting space is a finite direct sum of simple modules. Images of
a given simple module have zero projection to direct summands from
inequivalent simple modules, by Schur's lemma. Thus a finite
relation between sums belonging to pairwise inequivalent types has
each type contribution zero. The class reading distinguishes
different $\sigma$ by isomorphism class. Intersecting with the
level condition and using the injective target reading proves
independence of the native $\mathrm{isot}_j(\sigma)$.

Put $d_0=\dim E_{\rm base}$ for the reference base-level row space.
The member's dimension comparison equates
$\dim(W_{ip}\cap W_0)$ and $\dim(E_p\cap E_{\rm base})$ for every
$p$. The target and row direct sums and the isomorphism
$M_j\simeq W_0$ give $\dim M_j=d_0$. The base row reading identifies
$E_{\rm base}$ with the native space carrying the row basis. Consequently
$d_0=\sum_\sigma n_0(\sigma)$, where $n_0(\sigma)$ counts its basis
vectors of type $\sigma$. Independence now gives the exact remaining
test for the member's internal direct sum at row $j$:
\[
      \sum_\sigma\dim\mathrm{isot}_j(\sigma)=d_0 .
\]
It suffices to prove $\dim\mathrm{isot}_j(\sigma)\ge n_0(\sigma)$
for each type; summing and using independence forces equality. This
last inequality compares multiplicities of an individual
finite-adelic type across the two groups. More precisely, the upstream multiplicity input equates the cuspidal
multiplicity of an irreducible admissible good representation on
$Y_0$ with that of its transport by the finite-adelic isomorphism
to $W_i$, for the indicated archimedean discrete series. Once this
is supplied on the same member, the common finite level and good
projector act on the same finite representation, and taking the
level and projector subspaces preserves the multiplicity comparison.
The common spherical
label in L86M3 can contain multiple types and gives no such
lower bound by itself.
% END L86M5

\section{The two remaining compatibility conditions on the original member}

% BEGIN T86C1
\paragraph{T86C1. Recovering the parabolic from the character on a nonzero form.}
Fix a selected member, a coefficient index and a designated place above
one of the two type primes $g$. Put $n=r+t$, $k=\mathbb F_g$ and let
$U\subset\mathrm{GL}_n(k)$ consist of the block upper triangular
matrices having upper unitriangular first block of size $r$ and
arbitrary invertible second block. The principal and geometric frames
on this space give reductions related by $a\mapsto bab^{-1}$ for some
$b\in\mathrm{GL}_n(k)$. Indeed the standard frame in the member's
adelic realization identifies the action by conjugation over the local
field and carries the maximal compact into the stabilizer of the
lattice in the geometric frame. Every local integral invertible
matrix occurs, by inserting it at the split place with the
inverse-transpose partner and identity elsewhere. In the elementary
divisors of the frame change all exponents must be equal: otherwise
some elementary transvection would have a nonintegral conjugate.
Rescale the change of frame by a uniformizer power and reduce it.

If the row space is nonzero, $b$ preserves the first $r$ coordinates
as a subspace. To see this, take a nonzero realized form and a
linear functional on its value space detecting one of its values.
The resulting scalar function transforms through a character for
the inverse images of $U$ in both frames. In the principal frame
this is the specified level character, nontrivial on each simple
upper transvection in the first block by its additive character
on the sum of the superdiagonal. In the geometric frame this follows
by writing the same form using the member's adelic comparison:
the defining geometric level character acts by scalars on this
level subgroup. Matrices reducing to one fix the function: their
upper superdiagonal entries vanish in $k$, and the lower
determinant character on units is tame. These assertions only
use adeles supported at the tested place. Thus the stabilizer of
the line of this scalar function has a character on its reduction
$H$ containing $U$ and $bUb^{-1}$, nontrivial on each conjugated
upper simple root group $bU_{j,j+1}b^{-1}$, $1\leq j<r$.

Here is a short proof of the finite-group step. Write $b=b_1wb_2$
by Bruhat decomposition with $b_1,b_2$ upper triangular and $w$
a permutation. Upper triangular matrices normalize $U$. Conjugating
$H$ by $b_1^{-1}$ gives a group containing $U$ and $wUw^{-1}$.
The character on the latter is still nontrivial on each
$U_{w(j),w(j+1)}$: on the top block, upper triangular conjugation
scales each simple superdiagonal on the abelianization of the
upper unitriangular group. Its extra upper-right block contributes
no character (conjugation by a scalar on the lower block scales
that additive block by a scalar different from one). The bottom
diagonal block of a conjugated top transvection is the identity.
If either endpoint of such an edge were in the last $t$ coordinates,
a diagonal matrix in $U$ could scale this root group by $d\ne0,1$
by scaling just that lower coordinate. A character invariant under
this conjugation kills the whole group, since $(d-1)k=k$.
The type primes give $|k|>2$ and $r\ge2$, so these edges meet
every one of the first $r$ coordinates. The permutation preserves
their set, proving the assertion about $b$.
% END T86C1

% BEGIN T86C2
\paragraph{T86C2. The complex central action is the specified sum of right translations.}
At each designated place lift the first-block matrix of the finite
type group $G$ to an invertible integral matrix and adjoin the identity
lower block. Use the split unitary partner and take identity outside
the designated places. Denote these adelic lifts by $L(a)$, $a\in G$.
At auxiliary places their reduction is one. For a nonzero row, write
$A$ for the top block of the change $b$ in T86C1. In the geometric
frame $L(a)$ is block upper triangular with diagonal blocks
$AaA^{-1}$ and $1$. Conjugating the input $a$ by $A^{-1}$ at
each place gives the correct top coordinate simultaneously by a single
inner automorphism of $G$. The compensating level character in the
geometric action is one: the lower determinant blocks are identities,
the auxiliary reductions are one, and the top character of the
geometric flag level in this calculation is trivial.

For a rational conjugacy-class sum, the row's given integral-coordinate
comparison $\mathrm{AllT.cen}$ identifies the operator on rational
forms with its central level matrix. Under inclusion into ambient
coordinates that matrix acts as the sum of the corresponding finite
translations: the class sum commutes with the idempotent defining
the level. Under evaluation of an ambient coordinate form, the
translation is right translation by any lift multiplied by the inverse
level character of the lift, by the equivariance of the defining
function. The adelic comparison of the actual member identifies
these evaluated forms in the two frames. Use the lifts above and
reindex the class sum by the one inner automorphism. Thus for any
$z\in Z(\mathbb C[G])$ and form $F$ in this row,
\[
 \mathrm{real}(\mathrm{cenMt}(z)F)=
 \sum_{a\in G}z_a R(L(a))\,\mathrm{real}(F).
\]
Indeed the integral class sums span the complex center and the
rational coordinate basis spans the complex forms. For a zero row
the equality holds with the original lifts as well. This proves
$\mathrm{CenLoc}$ on the original localized module at every index.
The type space of each representative is the inverse image of its
isotypic sum under $\mathrm{real}$. Each right translation
preserves that sum. Thus every central operator preserves each type
space. The convention $\mathrm{onMt}=0$ away from the center
also gives the full stability input in M84T3. Together with
M84T3--4 and M84T7, this produces the integral action of the
center, its Hecke commutation and the scalar action on the whole
rational Hecke eigenspace at the original index.
% END T86C2

% BEGIN T86C3
\paragraph{T86C3. A local block vector and its two finite series.}
Fix an original representative $P$, including one with no Hecke
point, and a designated place $w$. Write $K=\mathrm{GL}_n(\mathbb
Z_g)$, $K_1=\ker(K\to\mathrm{GL}_n(k))$. Let $\pi$ be the
actual local factor of $P$ in the preferred split coordinate.
For $j=j(i)-1$ in zero-based notation let $\xi$ be the local
component of the tuple character $\mu_j^0\xi_{\rm tw}$, and
$\theta$ its character on $k^\times$. At either type prime
$\xi_{\rm tw}$ is unramified: the auxiliary character is supported
at an auxiliary prime distinct from both type primes. When
$\ell_i=g$ the remaining twisting character is trivial on units
above $\ell_i$; when $\ell_i\ne g$ the place belongs to $S_i$,
disjoint from its permitted ramification set $Q_i$. Hence
$\theta$ reads $\mu_j^0$ there. With the original residue
generator $h_1$, define $c\in k^\times$ by
$\theta(h_1)=\mathrm{Teich}(c)$ in the finite-series convention.
The tame-generator pin identifies $h_1$ with the residue of
$s(g^{1/(g-1)})/g^{1/(g-1)}$. A global integer unit at $w$
lifting this residue passes the exponent-one tame test as in
M84T10; the member's character test gives
$c\in\{\chi_J,\chi_J^{-1}\}$. The rank-one case of G85M2
with the same reciprocity convention identifies
$X-\mathrm{Teich}(c)$ with the inertia polynomial of
$\mathrm{locG}(\mu_C(i),w)$ at the named generator. Here the
actual character corresponds to the tuple by the member's
character identification; any unramified normalization is one on
inertia.

If a central top-block element sends a vector in $\mathrm{typeV}(i,P)$
to a nonzero vector, then there is $v\in\pi^{K_1}$ on which the
corresponding top action is nonzero, fixed by the parabolic
unipotent radical and transforming by $\theta\circ\det$ on the
lower block. Indeed in the level test the simple superdiagonal
character is trivial for lifts of this radical and of the lower
block, leaving respectively $1$ and $\theta\circ\det$.
T86C2 identifies the localized central operator with the finite
sum of top-block translations. For a lift of a single-place matrix
the other components may be taken to be identities; changing a
lift by the principal congruence group fixes the vector. Project
the resulting nonzero translated sum to a simple copy detecting
it in the global isotypic space, and then to a copy of $\pi$ in
the restriction to the local factor. Such copies span: translates
away from the place of an embedded local factor span a nonzero
global invariant subspace of the global irreducible. A sum of
simple modules is semisimple and the projections commute with the
local action. Project the original vector by the same maps.

In the series convention of the row, e.g. $\mathrm{grp.ser}$,
the idempotent for $f$ is the sum of the primitive central idempotents
of the finite irreducibles with Green label $f$. Finite parabolic
induction multiplies the cuspidal-support polynomials, by transitivity
of induction and uniqueness of finite cuspidal support. In rank $t$
the determinant character has label $(X-c)^t$: it occurs in
induction from the $t$ equal torus characters $\theta$, by
Frobenius reciprocity. The full inertia polynomial of any series
detected in $\pi^{K_1}$ is its lifted root polynomial at this same
generator. Apply the proof of G85M2 with ambient rank $n=r+t$.
Its only rank restriction was on an individual support factor $d$;
here $d\le n_V$ and $g-1$ is divisible by $Q_S q d_{\rm host}$
with $16\mid Q_S$, $q,d_{\rm host}\ge n_V$. Its numerical guard
therefore still holds. The member's $\mathrm{au}$, $\mathrm{embAV\_loc}$
and Weil--Deligne reading give precisely $\mathrm{locV}(i,P,w)$
in this computation, as in D85G1. These arguments supply the system
and all clauses of $\mathrm{LocType}$ in M84T8.
% END T86C3

% BEGIN T86C4
\paragraph{T86C4. The fourth tie.}
On the original member $D=w.D(Z)$, $\mathrm{SerLoc}(D)$ holds.
In the notation of T86C3 a top projector for $f$ detected on a
parameter space is detected in a representative's type space,
by the defining span for that parameter. Its projected vector is
fixed by the top projector, by the lower determinant projector,
and by the unipotent radical. Finite parabolic induction therefore
puts it in the series of $f(X)(X-c)^t$. On the member's base-change
list the rank-one string contributes $(X-\mathrm{Teich}(c))^t$;
all real shifts and the $\eta_T^{-1}$ twist preserve inertia.
Cancellation in the polynomial ring gives the lifted root polynomial
of $f$ on the remaining parameter. M84T8 proves equality with the
native monic label, including multiplicities, using the faithful
Teichm\"uller reading. T86C3 supplies its former hypotheses at
each original place and index. This is the original test for every
$F\in\mathrm{MtPar}$ and every polynomial $f$.
% END T86C4

% BEGIN T86C5
\paragraph{T86C5. Monodromy on the repeated determinant character.}
Suppose an actual finite irreducible top type $\tau$ occurs on
the space for $P$ in T86C3, with squarefree label $f$ avoiding
both prescribed values $\chi_J^{\pm1}$. A vector of the preceding
block kind can be projected to the $\tau$-isotypic space with
nonzero result. The lower block commutes with this projector and
the upper block normalizes the radical. Thus the parabolic
radical-invariants of $\pi^{K_1}$, restricted to its Levi, contain
$\tau\boxtimes (\theta\circ\det)$.
We show that the monodromy on the entire $n$-dimensional
parameter vanishes. The $t$ repeated copies of the same inertia
eigenvalue require an argument in addition to squarefreeness of $f$.

Write $A,B$ for the ranks $r,t$ finite representations
$\tau,\theta\circ\det$, and write $A*B$ for finite
Harish-Chandra induction. The finite product $A*B$ is irreducible;
moreover $A*B\simeq A'*B'$ with $A',B'$ irreducible in the same
respective cuspidal series forces $A'=A$, $B'=B$ up to
isomorphism. For clarity, Mackey's formula and Frobenius
reciprocity compute the intertwining space between these two
inductions by the double cosets of the $(r,t)$ parabolic.
Its intersections with a conjugate parabolic give the two
parabolic restrictions in the respective Levi blocks: on vectors
one takes invariants under the two intersecting unipotent radicals,
using their averaging idempotents in characteristic zero. In a cell
exchanging $a>0$ dimensions the exchanged $\mathrm{GL}_a(k)$
factor would be a constituent of a parabolic restriction from
each of the two series. Its cuspidal support would be a nonempty
submultiset of each original support, by transitivity and finite
Frobenius reciprocity. One support consists of $t$ copies of
$X-c$, while the other has no such root. Hence only the identity
cell contributes, giving
$\mathrm{Hom}(A,A')\otimes\mathrm{Hom}(B,B')$. Taking self
intertwiners and using finite semisimplicity proves irreducibility.
By the block vector and Frobenius reciprocity, $A*B$ occurs in
$\pi^{K_1}$.

Take the Zelevinsky multisegment of $\pi$. Its cuspidal factors
are depth zero, and their Green polynomials, including
multiplicities, have product $f(X)(X-c)^t$, by the support
calculation in G85M2 and T86C3. A factor with Green polynomial
$X-c$ is a rank-one character restricting to $\theta$ on units.
Indeed its Green polynomial has degree the cuspidal rank and
the degree-one Green correspondence reads the unit character.
Conversely every such character contributes $X-c$. Group the
segments on this unramified-twist line together, and put the
others in the first group. Segments in different groups are
unlinked. The multisegment classification and the irreducibility
of products on distinct cuspidal lines give
$\pi\simeq\pi_1\times\pi_2$ (normalized induction), with block
ranks $r,t$ and the second block on that line. Taking $K_1$
invariants gives finite parabolic induction of
$\pi_1^{K_{1,r}}\boxtimes\pi_2^{K_{1,t}}$ by Iwasawa
decomposition as in G85M2. All finite constituents of these two
spaces have respective labels $f$ and $(X-c)^t$ by the same
support argument. Decompose into finite irreducibles and apply
the finite intertwining calculation to the constituent $A*B$.
It gives $A\subset\pi_1^{K_{1,r}}$ and
$B\subset\pi_2^{K_{1,t}}$. In particular
$\pi_2\otimes(\xi^{-1}\circ\det)$ has a hyperspecial fixed
vector. Its spherical Weil--Deligne parameter has zero monodromy
by the unramified correspondence. Twisting back by a character
preserves this vanishing.

For $\pi_1$ the named inertia operator has characteristic
polynomial equal to the lift of $f$, with distinct roots.
Monodromy commutes with inertia, preserves every one-dimensional
eigenspace, and is nilpotent. Hence it is zero. The segment law
gives the direct sum of the two parameters for the irreducible
product. Thus $\mathrm{mono0}$ holds for the actual local
parameter, and its native label is $f(X)(X-c)^t$ by T86C3 and
the full-degree tame inertia reading of M84T9. Tameness here
follows either from the depth law applied to the $K_1$-fixed vector
as in D85G1 or from the individual depth-zero support factors.
At the conjugate place the paired split coordinates are exchanged
by inverse transpose. For irreducible admissible
$\mathrm{GL}_n$ this gives the contragredient by the
Gelfand--Kazhdan theorem. The duality and field-transport laws
of the correspondence apply: every support factor has
principal-congruence invariants as in D85G1. They give the dual
parameter under the conjugate field identification. Its monodromy
is the negative transpose of zero. M84T9 transports the designated
generator, so its label has degree $n$ and the inverses of the
roots of $f(X)(X-c)^t$, with multiplicity. The same $c$ works
in both conclusions, exactly as required in $\mathrm{admDZ}$.
% END T86C5

% BEGIN T86C6
\paragraph{T86C6. The integral witness, its parameter map and the second tie.}
Fix the local correspondence and the eventual design range of
M84T5. Choose for $\mathrm{lz}$ the finite irreducible
isomorphism classes, their ordinary character idempotents and
their Green labels in all ranks, with representations up to
isomorphism, ordinary occurrence, and finite parabolic induction.
In characteristic zero the finite group algebra is a product of
matrix algebras; these are all the primitive central idempotents,
orthogonal with sum one. The Green labels are monic of the
specified ranks with nonzero constant terms, and the member's
rank-$r$ series idempotents are the sums over their indicated
labels by T86C3. Take $\mathrm{kty}$ to be exactly the
irreducibles whose idempotents act nontrivially on some vector of
the given $\mathrm{typeV}$. Then T86C5 is the entire
$\mathrm{admDZ}$ assertion for this choice.

We supply the parameter-surjectivity field for the same witness.
If $\Sigma$ occurs at the original point $x$, the definition
supplies a representative $P$ over $x$ for which
$(\Sigma\sqcup\{(1,\mu_C(i),t)\})\otimes\eta_T^{-1}$
is a base change. Its parts have positive ranks and lengths,
are conjugate self-dual and are pairwise distinct as triples.
Twisting back is inverse on actual cuspidal representations
since the two characters multiply to one; this identity reaches
the original labels by the member's cuspidal identification
(equivalently by strong multiplicity one on the split-good
readings of the same global representations). Twisting by a
conjugate self-dual character preserves conjugate self-duality.
Deleting the final part and this common twist thus gives the
required positivity, self-duality and distinctness for $\Sigma$.
Its rank is $n-t=r$. At every designated type place the local
parameter of $P$ is tame: the local $K_1$-fixed-vector proof
in D85G1 uses the selected level at that place for every $i$,
irrespective of whether $g=\ell_i$. In the base-change sum
$\mathrm{isoLoc}(\Sigma,w)$ becomes a direct summand after an
unramified twist. Wild inertia therefore acts trivially on this
summand. In particular its monic label has degree $r$ at every
designated place by the wild-invariant label description in
M84T9. These are all the defining tests of the member's
concrete parameter set $\mathrm{ParC}$. Taking the inverse image
of this very list under the parameter identification
$\mathrm{par}.e_P$ gives one native parameter having the required
local values at \emph{every} place by its localization identity.
This proves $\mathrm{Par\_surj}$.

Finally take the roots of unity and coefficient-field series
coefficients of M84T5, the inertia characters and both polynomial
identities of M84T9, and the rational and integral actions and
projected eigenspaces of T86C2. The $\mathrm{EigSpan}$ assertion
is M84T3 on the same integral action. These choices give one
$Q\in\mathrm{Pend}(D)$ with both $\mathrm{admDZ}(Q)$ and
$\mathrm{EigSpan}(Q)$. In the eventual member quantifier the
finite list used in M84T5 was prescribed from the base before
the choice of $w,Z$. All later constructions are on the given
member and require no enlargement depending on $i,P$, a type
or a tested polynomial. Together with M84T10, M84T2 and T86C4
this proves all four original ties before making a good cut.
% END T86C6

\section{The inertia and monodromy conditions for the original rows}

% BEGIN I86A1
\paragraph{I86A1. The scalar on all inertia elements away from the coefficient prime.}
Fix a member, a coefficient index $i$ and a place $w\nmid\ell_i$.
Write $\xi_w$ for the local component of the tuple character at
$j=j(i)-1$ and $\lambda_w$ for its Weil character under the
member's local reciprocity convention. For every $s\in I_w$,
\[
                  \lambda_w(s)=\iota_i(\chi_1(i,s)).
\]
Indeed M84T10 identifies the global arithmetic line with the
$\ell_i$-adic avatar of the algebraic Hecke character obtained
from $\mu_C(i)\eta_T^{-1}$ with shift $-r/2$. Its proof compares
the specified split Frobenius values on a dense set; this part
uses only the coefficient index $i$, including when $\ell_i$
differs from both type primes. For an equivalent way to fix the
reciprocity convention, apply the rank-one Hecke-to-Galois
comparison to this algebraic character. The integrality of its
two infinity exponents follows from M84T6 and the computation
in M84T10. At a split good prime its Galois Frobenius value is
the unique Satake eigenvalue of that shifted character, exactly
the defining test for $\chi_1$. Density makes the two continuous
rank-one representations identical. Away from $\ell_i$ the
rank-one local comparison identifies the Weil character with
that of the shifted local Hecke character. The real shift and
$\eta_T^{-1}$ are unramified, and the native character reading
identifies $\mu_C(i)$ with the selected tuple. Restricting to
inertia gives the asserted equality with the original $\iota_i$
on the coefficient ring. Equivalently, in the idelic construction
the correction for the infinite exponents is supported above
$\ell_i$, so on units at $w$ only the same finite local
reciprocity value survives.
% END I86A1

% BEGIN I86A2
\paragraph{I86A2. Two spherical factors at the auxiliary places.}
Let $w\in S_i\cup Q_i$ satisfy the first disjunctive test in
$\mathrm{MonodromyInertiaType}$: either $w\in Q_i$, or
$\mathrm{Del}(i)$ holds and $w$ lies above $p_\delta$ or,
when $\ell_i=p_\delta$, above $p'_\delta$. Every such place
splits and is away from $\ell_i$. It is also away from both
type primes. For a $Q_i$ place this follows since all places over
a type prime away from $\ell_i$ belong to $S_i$ and $S_i,Q_i$
are disjoint. For the other disjunct it is the distinctness of
the auxiliary and type primes. Since $w\nmid\ell_i$, the
remaining $S_i$ case is over the active auxiliary prime, including
the switch to $p'_\delta$.

The unit character $\xi_w$ is tame and ramified here. In the
tuple $\mu_j^0\delta_{\rm act}\psi$ the first factor is
unramified away from the type primes. At a $Q_i$ place the
auxiliary factor is unramified by disjointness from $S_i$ and
$\psi$ is ramified by its exact $Q$-support. At a tested $S_i$
place $\psi$ is unramified and $\delta_{\rm act}$ has inertia
order 15. All these characters are tame on the indicated
units. Conjugate places carry inverse unit characters under
the CM identification, by the tuple's conjugate self-duality.
Write $k_w$ for the residue field. In a split coordinate supplied
by the level, the parahoric reducing to the $(r,t)$ parabolic
acts on any original level-character form by
$\theta(\det D)$, where $D$ is the lower diagonal block and
$\theta$ is the nontrivial character of $k_w^\times$ induced by
the tuple. Indeed the tested place belongs to the tuple's ramified
set by what we have just shown, so the member's $\mathrm{mt.para}$
applies at the place underneath it. Any local element of that
parahoric can be inserted as a single-place unitary adele. Its
principal congruence subgroup acts trivially. For every original
representative $P$ a nonzero vector in $\mathrm{typeV}(i,P)$
gives a nonzero vector with these transformations in its local
component $\pi$, by the isotypic projection argument of T86C3.

Here is the finite and local argument for any such residue field.
Its finite parabolic invariants contain
$1_{\mathrm{GL}_r(k_w)}\boxtimes(\theta\circ\det)$.
These finite factors have disjoint cuspidal supports, given by
$r$ trivial torus characters and $t$ copies of $\theta$.
The Mackey calculation in T86C5 applies over $k_w$, so their
finite parabolic induction occurs irreducibly in $\pi^{K_1}$.
In the multisegment support of $\pi$, each cuspidal factor has
principal congruence invariants, by exactness of compact
invariants and Iwasawa decomposition. For a depth-zero
supercuspidal, these invariants are its finite cuspidal type,
by depth-zero compact induction and the strict-gap Cartan
calculation in G85M2, valid for a local field with any residue
field. Finite parabolic induction of the support thus has the
same finite cuspidal support as its constituent just detected.
All factors here have rank one, on the two unramified-twist
lines of $1$ and $\xi_w$. This step uses no inertia-label or
large-residue-characteristic formula at the auxiliary prime.
Separate the lines by the unlinked-product theorem. The resulting
normalized product has blocks of ranks $r,t$. Exactness on the
maximal compact quotient and the identity-cell calculation in
T86C5 give the individual finite constituents $1$ and
$\theta\circ\det$ on those blocks. The first local representation
is therefore spherical and the second is spherical after twisting
by $\xi_w^{-1}\circ\det$. Both are supported on characters, so
the local correspondence's spherical, twist and product laws apply.
Their parameters have zero monodromy, the first is unramified
and the second acts on inertia through $t$ copies of $\lambda_w$.

If the chosen level coordinate corresponds to the other CM
place, use the inverse transpose between the two split coordinates
and tuple conjugate self-duality. By Gelfand--Kazhdan and the
duality and field-transport laws, as in T86C5, the dual parameters
give the two conclusions at $w$ with its own $\lambda_w$.
These laws apply here even for arbitrary residue degree since the
supports just obtained consist of characters. Transport by the
member's local reading and I86A1 give, for each $s\in I_w$,
\[
 \mathrm{mono0}(\mathrm{locV}(i,P,w)),\qquad
 \mathrm{inePol}(\mathrm{locV}(i,P,w),s)
       =(X-1)^r (X-\iota_i(\chi_1(i,s)))^t .
\]
% END I86A2

% BEGIN I86A3
\paragraph{I86A3. Extracting only the spherical block under the multiplicity guard.}
Let $w\in S_i$ lie above a type prime $g$. In particular
$w\nmid\ell_i$. Suppose the multiplicity of
$\alpha=\iota_i(\chi_1(i,\gamma))$ in the inertia polynomial
of $\mathrm{locV}(i,P,w)$ is at most $t$ at the designated
$\gamma=\mathrm{gam0}(i,w)$, with the native tame-generator
guard. We exhibit the two required local parameters at this
place. First work in the designated lift of the split pair.
The local component has a nonzero parabolic-radical fixed
vector of lower type $B=\theta\circ\det$, by the block argument
of T86C3 applied to the identity operator. Decompose under
$\mathrm{GL}_r(\mathbb F_g)$ and choose any irreducible top
type $A$ detected by this vector. For arbitrary $A$, finite
Frobenius reciprocity yields a common constituent of $A*B$ and
$\pi^{K_1}$. Its finite cuspidal support is the support of
$A$ plus $t$ copies of $\theta$. By the depth-zero support
calculation of G85M2, the number of factors on the local
unramified-twist line of $\xi$ equals the number of finite
cuspidal factors $\theta$ in this multiset. Each factor on that
line contributes one eigenvalue of $\lambda$ at \emph{every}
inertia element, since an unramified twist is trivial there;
the segment shifts preserve inertia as well. For the conjugate
coordinate the duality and transport argument of T86C5 gives
the same count with the conjugated inverse character, which is
$\lambda_w$ in that coordinate. All relevant support factors
are good here since the completion is $\mathbb Q_g$ and the
cuspidal factors have $K_1$-fixed vectors.

If $A$ contained even one extra $\theta$ in its cuspidal
support, the original inertia polynomial at $w$ would therefore
be divisible by $(X-\alpha)^{t+1}$, by I86A1. This contradicts
the assumed root multiplicity. Consequently the two finite
supports are disjoint and the second unramified-twist line has
exactly $t$ factors. The finite Mackey and unlinked-product
argument in T86C5 now gives a rank-$r$ first block and a rank-$t$
second block spherical after its character twist is removed.
No assertion about squarefreeness or monodromy of the first
block was required. Take the corresponding two Weil--Deligne
parameters, dualizing and transporting if $w$ is the conjugate
place. Their inverse images $a,b$ under the member's local
parameter identification satisfy
\[
 \begin{gathered}\mathrm{locV}(i,P,w)=a\oplus b,\quad
 \dim a=r,\quad\dim b=t,\quad N_b=0,\\
 \mathrm{inePol}(b,s)=(X-\iota_i(\chi_1(i,s)))^t\quad(s\in I_w).\end{gathered}
\]
In particular the proof uses the generator in the actual test;
equality with the Green-label generator is unnecessary.

Finally direct sums cancel at every place in the native parameter
space. Under its Weil--Deligne identification this is
Krull--Schmidt cancellation: the finite dimensional representations
with their monodromy operator form a category in which idempotents
split and each endomorphism has its invariant Fitting decomposition.
Finite dimension gives a finite decomposition into indecomposables,
whose endomorphism rings are local by Fitting's lemma, giving
uniqueness with multiplicities. Subtracting the multiplicities of
the common summand proves cancellation. I86A2, the displayed
decomposition and this unrestricted cancellation are the three
clauses of $\mathrm{At.MonodromyInertiaType}$. They prove the
original field $\mathrm{RowsB.a4}$ on every selected member.
% END I86A3

\section{Residual tame eigenvalues on an original member}

% BEGIN R86G1
\paragraph{R86G1. The cuspidal point and its inertia polynomial.}
Fix an original member $D=w.D(Z)$ with the layer-point guards, a location
$i$ of coefficient characteristic $\ell$, and a point parameter satisfying
the hypotheses of $\mathrm{g3p}$. Cuspidality makes its list
$\Sigma=[(r,\pi,1)]$. Its point supplies a unitary representative $P$ whose
base-change list contains $(r,\pi\otimes\eta^{-1},1)$. Here the cuspidal
embedding, its twist and localizations, and the characters are the ones
on $D$ itself, by B85U1 and the member's character frame.
In particular the automorphic representation $\Pi$ of $\mathrm{GL}_r$
represented by $\pi$ is cuspidal over the CM field $E$.

The base-change test makes $\Pi\otimes\eta^{-1}$ conjugate self-dual.
The restriction of $\eta$ to real-field ideles is the quadratic
character of $E/E^+$, so $\eta(x\bar x)=1$ by the norm law. Thus
$\eta^c=\eta^{-1}$ and untwisting gives $\Pi^c\simeq\Pi^\vee$.
At each embedding, the exponent $s$ of $\eta^{-1}$ is $1/2$ or
$-1/2$ by the infinity condition on the same frame. The shifted
infinitesimal multiset of $\Pi$, obtained by adding $s$ to each entry,
is a submultiset of the weight of $P$ by the infinity identity in
its base-change test. This last multiset has distinct entries in
$\mathbb Z+(n_V-1)/2$ by the algebraic coefficient weight.
Since $n_V=r+t$ and $s-t/2\in\mathbb Z$ for odd $t$, the
infinitesimal entries of $\Pi$ are distinct and in
$\mathbb Z+(r-1)/2$. This is regular algebraicity in the
normalization of B85U1.

Let $\rho$ be \emph{any} integral representation in the test.
Consider the characteristic-zero attachment $r_{\ell,\iota_i}(\Pi)$
with parameter normalized by $|\det|^{(1-r)/2}$ as in A72F8.
At a cofinite set of unramified places its characteristic polynomial
equals that of $\rho$ on the source's Frobenius elements.
Indeed $\mathrm{IsRepOfSig}$ reads the polynomial of the normalized
local parameter there. A72F7 and the good-support condition of the
cuspidal embedding identify the weak reading with classical local
Langlands. Both sides evaluate the same Weil element. When spherical
Satake coordinates are given at geometric Frobenius, evaluating at
arithmetic Frobenius inverts the entire list on both sides.
Chebotarev makes these Frobenius conjugacy classes dense. The
coefficients of both polynomials are continuous in
$\overline{\mathbb Q}_\ell$ after applying the specified coefficient
embedding to $\rho$, so they agree on every element of $G_E$.

At $w\nmid\ell$ the semisimplified local-global comparison suffices:
the Weil semisimplification of this attachment is the parameter of
$\Pi_w|\det|^{(1-r)/2}$ up to semisimplification
(I.~Varma, \emph{Local-global compatibility for regular algebraic
cuspidal automorphic representations when $\ell\ne p$},
Forum Math. Sigma 12 (2024), e21, Theorem~1).
On inertia the normalizing twist is trivial. Passing to
semisimplification preserves characteristic polynomials; the
unipotent exponential in the construction of the Weil--Deligne
representation also preserves them there, since it commutes with
the finite inertia matrix. A72F7 and B85U1 thus identify
$\operatorname{charpol}(\rho(u))$, in complex coefficients, with
$D.\mathrm{naut.inePol}(\mathrm{locG}(\pi,w),u)$ for every inertia
element $u$ at a designated type place $w\nmid\ell$.
% END R86G1

% BEGIN R86G2
\paragraph{R86G2. Simultaneous reduction at the designated generator.}
Put $g=g_k\ne\ell$, $w_+=\mathrm{plGw}(k,v)$,
$\gamma=\mathrm{tg}(k,v)=\mathrm{gam0}(i,w_+)$ and let $L$ be the
polynomial label of $\Sigma$ there. Write $S$ for its multiset of
roots in $\overline{\mathbb F}_g$. By full degree and the original
inertia-polynomial reading, R86G1 gives
$\prod_{a\in S}(X-\tau(a))$ for the polynomial on $\gamma$,
where $\tau$ is the member's multiplicative Teichm\"uller lift.
Every $a$ is nonzero since the matrix is invertible and
$\tau(0)=0$ (multiplicativity and injectivity on units force this).
Put $j=j(i)-1$. The member's separate test on its first line
$\chi_1$ gives $\chi_1(\gamma)=\tau(c)^{\delta}$ in the same
embedding, for $c=\chi_T(k,v,j)\in\mathbb F_g^\times$ and
$\delta\in\{1,-1\}$.

To reduce both formulas, embed $O_i$ in $\mathbb C$ as prescribed
and take its integral closure $B$ there. The roots of unity $\tau(a)$
belong to $B$. Choose a prime of $B$ over the maximal ideal of
$O_i$ by lying over. Its residue field is algebraic over the residue
field of $O_i$, so its embedding into $\overline{\mathbb F}_\ell$
can be chosen to extend the specified embedding $\mathrm{kFbar}_i$.
Composing $\tau$ with this reduction gives one multiplicative map
$J:\overline{\mathbb F}_g\to\overline{\mathbb F}_\ell$.
For $m$ prime to $\ell$, two distinct $m$th roots of unity have
different reductions: the factorization of $X^m-1$ in $B[X]$ would
otherwise have a repeated root after reduction, contrary to its
nonzero derivative at every root. Thus $J$ is injective on the
nonzero elements of order prime to $\ell$. Every $\ell$-power root
reduces to one. For a nonzero $a$ let $a_{\ell'}$ denote its
prime-to-$\ell$ component, the unique factor of order prime to
$\ell$ whose quotient from $a$ has $\ell$-power order. This is
precisely the CRT exponent in $\mathrm{primeToPart}$. Coefficientwise
reduction now gives
\[
 \operatorname{charpol}(\bar\rho(\gamma))
       =\prod_{a\in S}(X-J(a_{\ell'})),\qquad
 \bar\chi_1(\gamma)=J(c_{\ell'})^\delta.
\]
Here $\bar\rho=\mathrm{redF}(D,i,\rho)$ and $\bar\chi_1$ uses
the prescribed residue embedding. All multiplicities survive reduction.

By M84T9 the full degree of $L$ makes the local parameter trivial
on wild inertia, giving polynomial $(X-1)^r$ there in R86G1.
Each wild image of $\bar\rho$ is therefore unipotent, of
$\ell$-power order. The residual representation is continuous with
open kernel, and its restriction to the pro-$g$ wild group has
finite $g$-group image. Hence this image is trivial.
An arithmetic Frobenius lift $\phi$ acts by the $g$th power on tame
inertia, because the residue field at $w_+$ is $\mathbb F_g$.
Consequently
$\bar\rho(\phi\gamma\phi^{-1})=\bar\rho(\gamma)^g$ as matrices.
For the rank-one character the twist data express its value on each
wild element as a product of a $(g-1)$st root and an $\ell$-power
root. Its residual value has order prime to $g$, and is also in
the finite image of a pro-$g$ group, hence equals one. Thus the
product representation $(\bar\rho,\bar\chi_1)$ kills wild inertia.
% END R86G2

% BEGIN R86G3
\paragraph{R86G3. The same reduction for the two tests at a conjugate place.}
Every $w\mid g$ is either a designated $w_+$ or its CM conjugate,
as $g$ splits completely. At $w_+$ use R86G2 with signs
$(1,\delta)$. At its conjugate $w$ write $s=\mathrm{gam0}(i,w)$
and assume the required tame-generator condition there.
CM conjugation and an inner conjugation adjusting the prolongation
of the valuation transport $s$ to a tame generator $s'$ at
$w_+$. This uses transitivity on prolongations and the fact that
the field transport preserves inertia and its wild subgroup.

For clarity, the characteristic polynomial of $\rho(s)$ equals that
of $\rho((s')^{-1})$, and $\chi_1(s)=\chi_1((s')^{-1})$.
Indeed at cofinally many split good primes, conjugate
self-duality in R86G1 inverts the unitary Satake roots.
The common normalization multiplies their values at arithmetic
Frobenius by $q^{(1-r)/2}$. Thus conjugating $\rho$ has the
polynomials of its dual twisted by $\varepsilon_\ell^{1-r}$,
where $\varepsilon_\ell$ is cyclotomic. For $\chi_1$ the
Frobenius formula has power $(t-1)/2+(1-n_V)/2$; the two
conjugate values of each of $\mu_C$ and $\eta_T$ are mutual
inverses on the unramified split test. Therefore the product
of the two conjugate character values is
$\varepsilon_\ell^{t-n_V}$. These equalities extend by continuity
and Chebotarev: for each finite Galois quotient over $E$ take its
Galois closure over $E^+$ and prescribe a lift fixing $E$. A prime
with that Frobenius splits in $E$, and a choice of place realizes
the desired conjugacy class. At inertia away from $\ell$ the
cyclotomic character is one. Inner conjugation preserves
polynomials and character values, giving the asserted identities
also after reduction.

Write $A=(\bar\rho(\gamma),\bar\chi_1(\gamma))$, $M=\mathrm{ord}(A)$.
Both factors have open kernel and R86G2 kills wild inertia in their
product. The resulting finite image of inertia is cyclic.
Since $s'$ and $\gamma$ both generate its tame quotient
topologically before mapping, their images generate the same cyclic
group. The product value at $(s')^{-1}$ is $A^b$ for one integer
$b$ prime to $M$, including the trivial group of order one.
Let $\Theta$ be a multiplicative automorphism of
$\overline{\mathbb F}_g^\times$ acting by the $b$th power on
$M$-torsion. Explicitly decompose each torsion element into its
primary components. On its $p$-primary component raise to $b$ if
$p\mid M$, and act identically otherwise. These powers are
invertible on the specified components, so this gives the required
order-preserving automorphism. Extend it by $\Theta(0)=0$.

Each $J(a_{\ell'})$ is an eigenvalue of the first component of
$A$, and $J(c_{\ell'})^\delta$ is its second component's value.
Their orders divide $M$. Injectivity of $J$ on prime-to-$\ell$
torsion implies that the orders of $a_{\ell'}$ and $c_{\ell'}$
also divide $M$. Thus replacing $J$ by $J\circ\Theta$ raises
\emph{all} these values to the same $b$th power and preserves the
required injectivity. Triangularizing the matrix over the residue
algebraic closure shows that taking the $b$th power raises every
eigenvalue to $b$ with its multiplicity, regardless of Jordan
blocks. Hence the displayed formulas of R86G2 hold at $s$ with this
one new map and with the same $v,j$ and signs $(1,\delta)$.
Together with the Frobenius matrix equality of R86G2, this proves
both clauses of $\mathrm{ResidualTameEigenvalueCondition}$ for
every $i,\Sigma,\rho,k$ in the original test. It supplies
$\mathrm{GaloisPreconditions.g3p}$ on each selected member.
% END R86G3

% BEGIN R86G4
\paragraph{R86G4. Full degree on every point parameter and cuspidality on a good one.}
At a designated $w_+$ the local parameter of any $P\in\mathrm{AV}(i)$
is tame even when $g=\ell_i$. To see the latter case, apply the
fixed-vector argument of D85G1 directly. Its level input
$\mathrm{mt.uni}$ ranges over every coefficient place, every row and
each designated $k,v$; the tuple character is tame there also when
$g=\ell_i$, since both the reference character and the multiplier
are tame. A single-place principal-congruence element has trivial
level character. The nonzero type space of $P$ therefore yields a
nonzero locally fixed vector by the two projections of D85G1. The
depth law applies to its local representation, whose support is good
by the compact-averaging and parabolic-induction argument there.
This is entirely a statement on the complex parameter.

For any $\mathrm{HasParamX}$ witness $P$ the split identity expresses
this tame local parameter as the direct sum of the local strings of
$\Sigma$ and the distinguished string, with the common rank-one
twist $\eta^{-1}$. The member makes $\eta$ unramified at every
finite place. A direct summand of a representation on which wild
inertia acts trivially also has trivial wild action. Untwisting
therefore shows that the entire parameter of $\Sigma$ is tame.
Its label degree is the dimension of its wild invariants by M84T9,
and hence is its full dimension, for all $i,k,v$.
In particular this proves $\mathrm{GlobalDatumDepthZeroCondition}$
on the member for every list at a good point.

Suppose now that $\Sigma$ itself is good. Write its parts as
$(a_h,\pi_h,d_h)$. Positivity in the augmented base-change list
gives $a_h,d_h\ge1$ and $\sum_h a_hd_h=n_V-t=r$.
At each designated place, its label is the product of the
constituent labels to the powers $d_h$: local strings add direct
summands with unramified shifts and the inertia eigenvalues in
M84T9 multiply as polynomials under direct sum.
Every summand is tame, so each factor has degree $a_h>0$.
Squarefreeness forces $d_h=1$. If there were at least two
parts, the smallest $a_h$ would satisfy $1\le a_h\le r/2$,
and its factor would divide the label at \emph{every} designated
place above $g_0$. Goodness supplies a place where that divisor
degree is missing. Thus the list is $[(r,\pi,1)]$.
% END R86G4

% BEGIN R86G5
\paragraph{R86G5. The Hecke polynomial after removing the distinguished string.}
Choose one good point parameter as in R86G4, put
$K=\operatorname{Frac}O_i$, and take the attachment $R$ of $\Pi$
from R86G1 over a finite extension of $K$ in
$\overline{\mathbb Q}_\ell$. We construct the scalar twist explicitly.
The Hecke character $\alpha=\eta^{-1}|\cdot|^{-t/2}$ is algebraic:
the infinity exponents $\pm1/2-t/2$ are integers, since $t$ is odd.
Its $\ell$-adic avatar $\xi$ is unramified at every $w\nmid\ell$
because $\eta$ is unramified and the norm is trivial on finite
units. The construction is the one of M84T10: transport
$\alpha_f^{-1}$ to $\overline{\mathbb Q}_\ell$ and divide by the
embedding powers on the $\ell$-components to make the character
trivial on principal ideles; apply global reciprocity in the
arithmetic-Frobenius convention. For $q=\mathrm Nw$,
\[
       \iota_i\xi(\mathrm{Frob}_w)=\eta_w(\varpi_w)q^{-t/2}.
\]
These values belong to the image of $K$. Indeed the member's
column carries $\mathrm{ValLL0}$ on this frame. Its $\eta$ clause
places $\eta_w(\varpi_w)q^{-1/2}$ in the embedded small field,
and $q^{-(t-1)/2}$ is rational. The member's coefficient embedding
agrees there. Density of Frobenius and closedness of the finite
extension $K$ inside $\overline{\mathbb Q}_\ell$ give
$\xi(G_E)\subset K^\times$. Compactness makes its image under
the valuation a finite subgroup of $\mathbb Z$, hence zero.
Thus the original $O_i^\times$ contains all values of $\xi$.

Set $R'=R\otimes\xi$ and choose over a finite coefficient extension
$K'/K$ a stable lattice. Such a lattice is the span over $O_{K'}$
of all translates of a basis lattice: compactness bounds them
between two scalar multiples of that lattice. On every split
place $w$ in the $\mathrm{SplitOut}$ test the member's $S_i,Q_i$
exclusions imply that $w$ lies above neither a type prime nor the
selected auxiliary prime, and $w\nmid\ell$. The local parameter of
$P$ is unramified by its native unramifiedness clause.
Its base-change identity therefore makes the constituent
$\Pi_w\eta_w^{-1}$ unramified, including zero monodromy.
The classical comparison for the polarized cusp in A72F8 gives
unramifiedness of $R$ there and the normalized Satake polynomial.
The point equality $\mathrm{pt\_loc}$, evaluated at arithmetic
Frobenius $F$, gives
\[
  H_{x,w}(X)=\operatorname{charpol}(R'(F))
                  \prod_{c=1}^{t}(X-\chi_c(F))
\]
in the common coefficient field. In fact the $\Pi$ block in the
augmented list acquires precisely the above scalar, since
$(1-n_V)/2=(1-r)/2-t/2$. The first distinguished root is
$\chi_1(F)$ by the character normalization, and the $k$th
successive shift multiplies it by $q^{-k}$, exactly
$\chi_{k+1}(F)$ for $0\le k<t$. This verifies every factor.
Monic division on $O_i[X]$ shows that the polynomial of $R'(F)$
already has coefficients in $O_i$. Split Frobenius density as in
R86G3 and closedness of $O_i$ give the same assertion for all
$g\in G_E$.

For every residual representation $r_b$ in $\mathrm{IsRhoBar}$,
its defining polynomial times the same distinguished product is
the reduction of $H_{x,w}$. Cancel that monic product and use the
same density, now with locally constant polynomials, to obtain
\[
 \operatorname{charpol}(r_b(g))=
       \operatorname{charpol}(\overline{R'}(g))\quad(g\in G_E)
\]
in $\overline{\mathbb F}_\ell$, using a residue embedding for
$K'$ extending the one for $O_i$.
% END R86G5

% BEGIN R86G6
\paragraph{R86G6. Missing divisor degrees imply the original residual irreducibility test.}
Let $U$ be either the reduction of $R'$ on its stable lattice
or any $r_b$ as in R86G5, extended to $\overline{\mathbb F}_\ell$.
The reduction of $R'$ at inertia away from $\ell$ has exactly the
polynomials of $R$ there since $\xi$ is unramified. Apply the
proofs of R86G1--3 to this stable lattice over $O_{K'}$, using the
extended residue embedding. The designated polynomial and the
two conjugate-place formulas therefore hold also for $U$ by
R86G5. For the designated Frobenius \emph{matrix} identity on an
arbitrary $r_b$, note that its wild inertia polynomial is
$(X-1)^r$ by the same equality. Its kernel is open by assumption.
Its wild image is consequently trivial by the coprime-orders
argument of R86G2, which gives the matrix identity on $r_b$ itself.

The two bad sets $\mathrm{BB}(r,g_k)$ are disjoint in the
original type data. Choose $k$ with $\ell$ outside its bad set,
so $\ell\ne g_k$. At each $v$ above this $g_k$, all eigenvalues
of $U(\gamma_v)$ are distinct: squarefreeness gives $r$ distinct
roots $a$, the good tuple separates their $a_{\ell'}$, and $J$ is
injective on those. If $U$ had a nonzero proper invariant subspace,
take that space or its quotient of dimension $d\le r/2$.
Its $\gamma_v$-eigenvalues form a $d$-element subset of the
full list. The Frobenius matrix identity makes this subset stable
under $g_k$th powers. The bijection $a\mapsto J(a_{\ell'})$ on
the label roots commutes with $g_k$th powers; projection to
prime-to-$\ell$ torsion and $J$ are both multiplicative there.
Consequently the preimage gives a Frobenius-stable $d$-element
subset of the label roots, hence a degree-$d$ divisor over
$\mathbb F_{g_k}$. This holds for every $v$, contrary to goodness
for that degree. This proves absolute irreducibility of $U$.

For \emph{each} $k$ with $g_k\ne\ell$, regardless of its bad set,
the original $S_i$ definition puts every $w\mid g_k$ in $S_i$,
so $\mathrm{gam0}(i,w)$ satisfies the tame-generator test.
The character and matrix polynomials of R86G2--3 use one map
injective on prime-to-$\ell$ torsion. If the first character
were an eigenvalue, the prime-to-$\ell$ part of a label root
would equal that of $\chi_T(k,v,j)$ or its inverse. The two
goodness tests exclude both at every such $\ell$.
Together these results prove $\mathrm{GaloisPreconditions.rhoBar}$:
$\mathrm{HG}$ supplies the good point and list used above, and
the argument covers every residual test representation over its
original residue field. Root exclusion descends because that
field embeds in $\overline{\mathbb F}_\ell$.
% END R86G6

% BEGIN R86G7
\paragraph{R86G7. Integral descent at a good point.}
For each good point $x$ and list $\Sigma$ in the integral residual representation condition,
R86G4--6 apply to this very list. We descend $R'$ from $K'$ to
the specified $O=O_i$. In a basis of the stable $O_{K'}$-lattice,
write $B$ for the $O$-span of its group matrices. This is a
unital algebra because the product of two group matrices is a
group matrix. Put $k=O/\mathfrak m$, $k'=O_{K'}/\mathfrak m'$.
Absolute irreducibility in R86G6 and Burnside's matrix algebra
theorem over $\overline{\mathbb F}_\ell$ give $r^2$ group
matrices $b_1,\ldots,b_{r^2}$ whose reductions form a basis of
$M_r(k')$. In particular they form a basis over $K'$. The trace
pairing on the full matrix algebra is perfect in every
characteristic (dual matrix units are transposes). Hence
$\det(\operatorname{tr}(b_j b_h))$ is a unit in $O_{K'}$.
All entries lie in $O$ by R86G5 applied to products, so this
determinant is a unit in $O$. The coordinates of every group
matrix in the $b_j$ basis are obtained by inverting this Gram
matrix against its traces of products with the $b_h$. They
belong to $O$, proving $B=\bigoplus_j O b_j$.

Reduction and extension to $k'$ now identify
$(B/\mathfrak m B)\otimes_k k'$ with $M_r(k')$. Thus
$B/\mathfrak m B$ is central simple over $k$ (simplicity and
the center can both be tested after this field extension).
The residue field $k$ is finite, so its central simple algebra
is split by the finite division-ring theorem. Lift a rank-one
idempotent of $M_r(k)\simeq B/\mathfrak m B$ to an idempotent
$e\in B$. For example, start from any lift $a$ and iterate
$a\mapsto 3a^2-2a^3$. This fixes $a$ modulo $a^2-a$ and
squares the ideal of the error. The sequence converges in $B$,
which is finite free over the complete DVR $O$. Its limit works
also in residue characteristic two.

The left ideal $Be$ is a direct summand of a free module over
$O$; reduction identifies it with a column module of dimension
$r$, so it is free of rank $r$. The action map
$B\to\operatorname{End}_O(Be)$ is an isomorphism by reduction
and Nakayama (both free ranks are $r^2$). After extending to
$K'$ its simple module of dimension $r$ is isomorphic to the
original matrix representation, since a full matrix algebra
has one simple module up to isomorphism. We have thus obtained
$R'_O:G_E\to\mathrm{GL}_r(O)$ with the same characteristic
polynomials. Its continuity follows from the Gram-matrix
coordinate formula and continuity of the original matrices:
the coordinates lie in the closed subfield $K$ with its induced
topology. Its generic fiber is absolutely irreducible by the
full matrix algebra just computed.

Set $\rho=R'_O\otimes\xi^{-1}$ with $\xi$ from R86G5.
This gives $\mathrm{IsRepOfSig}$ on the original lattice, since
over $K'$ it is the ordinary attachment $R$. The twist $\xi$
is trivial on inertia at every type prime away from $\ell$,
and $\rho\otimes\xi=R'_O$ satisfies the Hecke polynomial
identity at every $\mathrm{SplitOut}$ place by R86G5.
These are exactly the three clauses of the original
$\mathrm{GlobalDatumResidualRepCondition}$, for the specified
$x$ and $\Sigma$. Thus the $\mathrm{rho}$ and $\mathrm{dz}$
inputs of the row conditions also follow on original members.
% END R86G7

\section{Recovering the lower parameter from an upper theta lift}

% BEGIN U86T1
\paragraph{U86T1. The split parameter of a nonzero pairing.}
Write $T$ for the up operator at the original localization $i$.
Suppose $\langle F,Tf\rangle\ne0$ for some $F$ of upper type $P$
and some $f$ in the lower level space of type $\sigma$. Fix any
split finite place $w$ of $E$ and put $n=r+t$. In the two general
linear groups let the components be $\sigma_w$ and $\pi_w$.
Then, in the member's local-parameter convention,
\[
 \operatorname{rec}(\pi_w)=
 \left(\operatorname{rec}(\sigma_w)\oplus
      \bigoplus_{k=0}^{t-1}\operatorname{rec}(\mu_{i,w})
                  |\cdot|^{(t-1)/2-k}\right)\otimes
                     \operatorname{rec}(\eta_w^{-1}).
 \tag{1}
\]

Realize $F$ as a definite algebraic modular form and $f$ in the
cuspidal multiplicity space with its invariant inner product.
On the upper side the member's adelic identification gives the
weight intertwiner $e_W$; its inverse identifies the weight values
of these forms. The rational group actions agree in the two standard
frames, and the member's evaluation identity transports the functions
by this same intertwiner. On the lower side the two readings give
isomorphic copies of the same finite type by the row-class
identification. The coefficient and tuple identifications are those
of M84T1 and M84T6. Thus the principal-level map in L86M1 evaluates,
up to the common nonzero metric normalization, by the global theta
integral on these forms and the prescribed Schwartz vector.
The adjoint identity $\langle F,Tf\rangle=\langle T^*F,f\rangle$
puts the lower form in the conjugate-linear variable of the
trilinear pairing (or take the complex conjugate of the entire
identity for the opposite inner-product convention).
The integral is invariant under simultaneous translation of the
kernel and forms by change of the two integration variables.
The upper quotient is compact and cuspidal rapid decay on the
lower quotient justifies this change for the theta kernel.

At $w$ the Schwartz space is spanned by local tensors with
other coordinates fixed. For each finite-adelic irreducible,
fixing the other tensor coordinates likewise gives equivariant
copies of its local component; these span the smooth vectors.
For an isotypic space take the sum of the images of the irreducible.
Successively expand the three arguments in these spans. Some triple
has nonzero value. With the other coordinates fixed the integral is
a nonzero local trilinear pairing. Translations supported at the split
place realize arbitrary matrices in both factors. In the split
Schwartz model the upper determinant contributes the scalar
$\eta_w(\det h)$; putting it on the upper vector gives
$\pi_w\otimes(\eta_w\circ\det)$. In the lower variable use its invariant Hermitian form to
identify the antilinear functional with a vector in the smooth
lower representation. The linear upper variable gives a vector in
the smooth contragredient of the twisted upper representation.
For a fixed smooth Schwartz vector its pairing is fixed by an
open compact product. Averaging each automorphic factor over its
compact subgroup factors the functional through their two spaces
of fixed vectors. Those are finite-dimensional by admissibility,
so their product dual is the tensor of their duals. Varying the
Schwartz vector gives the equivariant quotient map required by
the type-II correspondence. This gives the type-II pairing for $\sigma_w$ with splitting character
$\mu_{i,w}$.

In the explicit type-II correspondence, passing from rank $r$ to
$r+t$ inserts $t$ consecutive exponents centered at zero in the
contragredient Langlands data (A. M\'inguez, \emph{Correspondance
de Howe explicite}, Ann. Sci. ENS 41 (2008), Theorem 1).
The upper factor in the quotient is the smooth contragredient; taking its dual cancels the contragredient of the lower Langlands data.
In the rank-$(r,r+t)$ matrix model the character multiplier on
$(g,h)$ is $\mu_{i,w}(\det g)\mu_{i,w}(\det h)^{-1}$.
The determinant characters in this action replace the trivial
inserted block by $\mu_{i,w}\circ\det$: twist the lower input by
$\mu_{i,w}^{-1}$ before applying the formula and twist the inducing
factors back by $\mu_{i,w}$. The lower representation is
unitarizable: pull the positive invariant form back to a nonzero
copy of the finite type, then to the embedded local component.
The normalized product
$\sigma_w\times(\mu_{i,w}\circ\det)$ is irreducible by Bernstein's
unitary product theorem. Its parameter is the direct sum of the
parameters of its factors; that of the determinant character
consists of the displayed $t$ lines with zero monodromy. For the
upper label, H86P1 and B85U1--2 give good split supercuspidal
support at every split place: the original upper label is
spherical at a nonsplit place, and the member's reading of this
predicate includes the support condition. Twisting preserves
good support. In the irreducible product the two inducing supports
form the support of the product, so the lower support is good as
well. Thus A72F7 identifies both weak parameter readings with
their classical ones. Reading the resulting parameter and undoing the upper
twist proves (1). In particular, if the entire nonzero image
$T(\mathrm{isot}_j(\sigma))$ is in type $P$, take $F=Tf\ne0$.
This proves the original $\mathrm{adm3}$.
% END U86T1

% BEGIN U86T2
\paragraph{U86T2. The row bound and uniqueness of the upper label.}
At a split place where its parameter is unramified, the component
of a nonzero lower row is spherical by the spherical local
correspondence: the classical inverse of an unramified parameter is
spherical and has character support, on which the weak and classical
maps agree. The weak-map injectivity with one good input in A72F7
therefore identifies the lower component. It is unitary by the invariant-form argument of
U86T1, which only uses occurrence of the lower type. Write
$q=\mathrm Nw$. Each spherical unitary Satake root, and its
inverse, has absolute value at most $q^{(r-1)/2}$. Indeed, in the
spherical case of Tadi\'c's unitary dual classification a unitary
character block of length $b$ has exponents from $-(b-1)/2$ to
$(b-1)/2$. A complementary pair of such blocks adds opposite
shifts of absolute value less than $1/2$ and consumes $2b$ of the
rank. Both kinds give the required bound. Only unramified character
blocks occur in a spherical product, by the hyperspecial invariants
in normalized induction and the spherical classification. Since
$q>1$ and $r<t$, this gives the strict bound at exponent $(t-1)/2$,
also for the inverse-Satake reading. This is $\mathrm{RowsB.rowSize}$.

Suppose $P,P'$ both pair nontrivially with this lower row.
Equation (1) gives equal split parameters at every split place.
Take their base-change lists from H86P1 and the representations
of H86P2. At good split Frobenius elements these have equal
characteristic polynomials. By S86A1 the semisimple representations
are isomorphic. H86P2 now gives equal list parameters at almost
all places, including nonsplit ones, after undoing the common
unramified twist. S86A1 recovers equality of the lists up to
permutation. At every nonsplit place both labels are spherical
and have equal unramified base-change multisets by H86P1. Remove
one $1$ in odd rank. On each two-element inversion orbit the number
of pairs is its common multiplicity; at $1,-1$ the remaining
multiplicities must be halved. This determines the relative tuple
up to permutations and individual inversions, precisely the
relative Weyl action. Thus the nonsplit representations agree.
At split places use (1) and injectivity of the local correspondence
on the supported components. Taking the restricted tensor product
of these component isomorphisms, normalized on spherical vectors
almost everywhere, identifies the finite-adelic irreducibles.
The upper labels enumerate their isomorphism classes in fixed
weight, so $P=P'$ even if the representation has multiplicity.
% END U86T2

% BEGIN U86T3
\paragraph{U86T3. The whole lower image lies in a single upper type.}
At fixed $i$ the upper level space is spanned by its type spaces.
They are orthogonal in the native inner product: the member reads
this product through the definite invariant class-sum pairing.
A nonzero cross pairing between two simple copies would give a
nonzero intertwiner from one to the Hermitian dual of the other.
That dual is isomorphic to the other unitary representation, so
Schur's lemma gives orthogonality for inequivalent copies. Taking
sums and imposing the level preserves this orthogonality.

If the lower summand is nonzero, L86M4 gives $T^*Tf=\lambda f$
on it with $\lambda>0$. Choose $f\ne0$. Since
$\|Tf\|^2=\lambda\|f\|^2>0$, some upper type $P$ pairs nontrivially.
For any $f'$ in the same lower summand every component of $Tf'$
outside $P$ is zero: a nonzero orthogonal projection would provide
a second nonzero pairing, contradicting U86T2. Thus $P$ satisfies
$\mathrm{thetaOf}(i,\sigma,P)$. The nonzero image and disjointness
of the type spaces make $P$ unique. This proves
$\mathrm{RowsB.thetaOne}$. With U86T1--2 and I86A3 it gives all
four fields of $\mathrm{RowsB}$.
% END U86T3

% BEGIN U86T4
\paragraph{U86T4. The long string, the good rows and the original first conclusion.}
Fix this $P$ and a base-change list $\mathcal P$ for it.
At almost every split place all components are spherical. In (1)
there is exactly one root of absolute value $q^{(t-1)/2}$, the
extreme root from the distinguished character string. All lower
roots are smaller by U86T2, and no root of (1) is larger.
Every unramified root of a unitary cuspidal base has absolute value
strictly less than $q^{1/2}$ by S86A2. Hence a block of length
$d<t$ in $\mathcal P$ has all roots strictly smaller than
$q^{(t-1)/2}$. There must be a block of length $d\ge t$.
The total rank is $r+t<2t$, so exactly one block has this property
and its cuspidal base is a character. Unitarity makes its extreme
absolute value exactly $q^{(d-1)/2}$, whence $d=t$. Equating the
two extreme roots at every such split place and cancelling the
common positive norm factor identifies this character locally with
$\mu_i\eta^{-1}$. The rank-one comparison in S86A2 identifies it
globally, on the member's original labels.

Remove this block and twist each remaining cuspidal base by $\eta$
to form $\Sigma$. It has positive parts and rank $r$. At every
split place untwist (1) and cancel the common character string
using I86A3, direct-sum cancellation for Weil--Deligne parameters.
This gives precisely $\mathrm{HasParamW}(i,\sigma,\Sigma)$,
including ramified split places. Appending $(1,\mu_i,t)$ and
twisting by $\eta^{-1}$ recovers $\mathcal P$ up to permutation,
so all six base-change identities, including the nonsplit and
infinity identities, hold for this reconstructed list. U86T3
makes the assertion valid for every $P'$ with $\mathrm{thetaOf}$.

Finally suppose the row is good and let $\Sigma'$ be \emph{any}
list satisfying $\mathrm{HasParamW}$ for it. At each designated
place its local reading equals the lower reading. The row label is
the polynomial of the lower reading on wild invariants by the
member's finite-type comparison (M84T9 and G85M3). Goodness gives
full degree $r$, and the sum of the ranks of the parts, counted
with lengths, is $r$. Each factor must have its full positive
degree. A length $d>1$ contributes the $d$th power of that factor,
contrary to squarefreeness since norm shifts leave inertia unchanged.
If several length-one blocks remain, one has degree $a$ with
$1\le a\le r/2$. It would give a divisor over the residue field
of degree $a$ at every designated place over the first type prime.
The missing-divisor-degree clause supplies a place without such
a divisor. Thus $\Sigma'$ consists of one rank-$r$ cusp of length
one. This proof also covers lists for which conjugate self-duality
was not assumed.

Together with L86M4 and U86T3 this proves all three assertions of
$\mathrm{Conclusion42}$. H86P4 proves $\mathrm{provisionFortyOne}$
and S86A4 proves the unused middle conjunct. Hence the original
member field $\mathrm{GaloisPreconditions.Conclusion1}$ holds.
% END U86T4

\section{The residual representation of a row and its upward theta image}

% BEGIN J86O1
\paragraph{J86O1. Finding the Hecke point at the same coefficient place.}
Fix an original member, a location $i$ and a nonzero row
$M_j[\sigma]$, with $j=j(i)$ and $\lambda=w(i)$.
Take its theta type $P$ from U86T4. Its image is nonzero and
lies in $V_P=\mathrm{typeV}(i,P)$. We need a point for this
lift at some location with the same pair $(j,\lambda)$.

Use the original $\mathrm{NdSpec}$ for this pair. Write $O_c$ for
its common coefficient DVR, $K_c$ for its fraction field, and $T_c$
for the $O_c$-algebra generated by the specified split spherical
operators on the full rational Hecke space.
By M84T7 each generator acts on $V_P$ by a scalar, the same for all
the allowed frames and uniformizers in its test. The rational basis
of forms gives both the complex basis and the $K_c$-basis by the
member's two basis identifications. In those bases each spherical
operator is the image of the same rational matrix, by
$h_L$ and $h_K$ in the Hecke model. The embedding of $K_c$ in
$\mathbb C$ is the extension of $\mathrm{Sp.iotaC}$, compatible
with those bases by $\mathrm{phK\_emb}$ and
$\mathrm{iotaC\_coh}$. Each such matrix splits over $K_c$ by the
full generalized eigenspace decomposition. Its eigenvalues belong to
$O_c$ by $\mathrm{Sp.std}$. In particular all scalars on $V_P$
are images of elements of $O_c$.

Choose one nonzero vector in $V_P$. The simultaneous equations for
these eigenvalues over $K_c$ have a nonzero solution there: their
linear functionals span a finite-dimensional dual subspace, a finite
basis of which has the same matrix rank after extension to
$\mathbb C$. Acting on a nonzero solution yields a character
$\psi:T_c\to O_c$. Indeed, each element is an $O_c$-polynomial
in the operators, hence acts by its evaluation at the integral
scalars. A relation among the operators kills this vector, so that
its scalar evaluation vanishes. The character fixes $O_c$ and its
simultaneous eigenspace is nonzero. Its residue kernel is maximal,
since the constants surject onto the residue field.
The member has $\mu_{\rm At}(\lambda)\ne\mu_0$ at every actual
location. Thus $\mathrm{Sp.mx\_all}$ gives a location $i'$ of
the same row and place with this maximal ideal, and
$\mathrm{Sp.ptX}$ gives $x'\in X_{i'}$ for $\psi$.

The tuple at both locations is $\mu_{\rm Tw}(\lambda)$; the
representation occurs in the same automorphic space with the same
level. By surjectivity onto its occurring irreducible types, choose
$P'$ at $i'$ representing this same type. The two type subspaces are
identified, as are the original theta operators on $M_j$. Thus
$P'$ is a theta type of $\sigma$. Its generators act by the
values of $x'$ by $\mathrm{Sp.ptX\_hT}$ and the common coefficient
embedding, on every vector of the type. The point test of M84T7
gives $\mathrm{ptV}(i',P')=x'$.
% END J86O1

% BEGIN J86O2
\paragraph{J86O2. Joint occurrence with the prescribed residue map.}
Take any $k,v,i,\sigma$ in the original joint occurrence test,
with $g=g_k\ne\ell_i$, and put $w=\mathrm{plGw}(k,v)$,
$\gamma=\mathrm{tg}(k,v)$. Use $i',P',x'$ from J86O1.
By U86T4 choose a row list $\Sigma$ at $i'$ whose augmentation
is a base-change list of $P'$. Goodness of $\sigma$ and the
nonzero row make it a singleton cusp list. At all designated
places the list has exactly the row labels, since these places
split and $\mathrm{lab\_loc}$ reads their labels. It follows that
$\Sigma$ satisfies both goodness tests with the original characters.
The distinguished term in its augmentation puts $x'$ in $Z$,
and the same list witnesses $x'\in Z_{\rm good}$. Apply R86G7 at
this point to obtain the \emph{untwisted} integral representation
$\rho$ of $\Sigma$ over $O_{i'}$.

The original datum has the identity $D.\mathrm{rbar\_tie}$:
for every list satisfying $\mathrm{HasParamW}$ and every integral
$\mathrm{IsRepOfSig}$ representation, it identifies the
characteristic polynomial of $D.\mathrm{ntypes.rbar}(\sigma,w(i'))$
at \emph{every} group element with the reduction of that integral
polynomial through $\mathrm{kFbar}_{i'}$. This applies to the pair
just constructed. Both $w(i')$ and $w(i)$ equal $\lambda$, so the
residual matrices in the required test are the very same matrices.
The identification of residual characteristics is the canonical one
from $w(i')=\lambda$, $\ell_{i'}=\operatorname{resChar}(\lambda)$.

For the reduction of the roots in R86G2 take the member's own
valuation ring $B=A.\mathrm{pl.plB}(\lambda)$ in $\mathbb C$ and
the residue identification $\kappa=A.\mathrm{pl.pl}\kappa(\lambda)$.
The equality $\mathrm{Au.kFbar}$ makes $\kappa$ after reduction
on $O_{i'}$ the specified map on integral coefficients.
The equality $A.\mathrm{iota}$ makes it on each Teichm\"uller
value $\tau(a)$ exactly $D.\mathrm{ntypes.iota}(k,v,\lambda)(a)$.
All these roots of unity are units of $B$. Reduction kills their
$\ell_i$-power factors. Thus with $L=D.\mathrm{nd.lab}(k,\sigma,v)$
and $a_{\ell_i'}$ the prime-to-$\ell_i$ part, R86G1--2 give
\[
 \det\!\left(X-\bar r_{\sigma,\lambda}(\gamma)\right)
   =\prod_{a\in\operatorname{roots}(L)}
        \left(X-D.\mathrm{ntypes.iota}(k,v,\lambda)(a_{\ell_i'})\right),
\]
including all multiplicities in the original field.

The wild inertia polynomial of $\rho$ is $(X-1)^r$ by the full
degree tame reading in R86G1--2. The same polynomial holds for the
actual $\bar r_{\sigma,\lambda}$ by $D.\mathrm{rbar\_tie}$. The
native residual representation has open kernel, hence its image of
wild inertia is a finite $g$-group. Each of its matrices is unipotent
in characteristic $\ell_i$ and has $\ell_i$-power order. The two
orders are coprime, so this image is trivial. Arithmetic Frobenius at
the designated place acts on tame inertia by the $g$th power since
the completion is $\mathbb Q_g$. A lift $\phi$ in the decomposition
group consequently satisfies
$\bar r_{\sigma,\lambda}(\phi\gamma\phi^{-1})
=\bar r_{\sigma,\lambda}(\gamma)^g$. This proves
$\mathrm{RowsC.c4b}$ on the original member.
% END J86O2

% BEGIN J86O3
\paragraph{J86O3. The upward inclusion at each specified point.}
Fix $i$, $x\in Z_{\rm good}$ and any $\Sigma$ satisfying
$\mathrm{HasParamX}(i,x,\Sigma)$. Choose a witness $Q$ with
point $x$ and with the prescribed augmented list. It suffices to
test a summand $M_j[\sigma]$ of $M_{j,\Sigma}$, so assume
$\mathrm{HasParamW}(i,\sigma,\Sigma)$ and the summand is nonzero.
By U86T4 its upward image lies in some $V_P$ and the lift $P$
has the augmented base change of a row list $\Sigma'$.
At each split place the two row lists $\Sigma,\Sigma'$ read
the same parameter $\mathrm{locW}(i,\sigma,w)$. Appending the
same distinguished string and tensoring by $\eta^{-1}$ give
the same local parameter, by the direct-sum and tensor-product
formulas. Hence $\mathrm{locV}(i,P,w)=\mathrm{locV}(i,Q,w)$
at each split place.

At each place in the split-good test the hyperspecial operator
of index $b$ acts by a scalar on $V_P$ by M84T7. The member's
spherical polynomial formula identifies $(-1)^b(\mathrm Nw)^{b(b-1)/2}$
times this scalar with the coefficient of $X^{n_V-b}$ of the
normalized local polynomial of $P$. Indeed one may evaluate on a
nonzero fixed vector in the irreducible representation of $P$:
project a nonzero level vector in $V_P$ onto one of its irreducible
copies and intertwine this copy with the chosen representative.
The projected vector remains hyperspecial fixed and has the original
Hecke scalars, since the operators are sums of translations. This
proves the coefficient identity for every frame and uniformizer in
the point test. The polynomial just identified equals the one for
$Q$, which is the Hecke polynomial of $x$ by $\mathrm{pt\_loc}$.
Each coefficient factor is nonzero, so all tested scalars on $V_P$
are the specified ones for $x$. The member's converse point test
gives $\mathrm{ptV}(i,P)=x$. Consequently the upward image of this
summand lies in $V_P\subseteq\mathrm{eigC}(i,x)$. Summing over
$\sigma$ proves $\mathrm{RowsC.parUp}$ for every tested $x,\Sigma$.
% END J86O3

\section{A good signature for each prescribed lower-row parameter}

% BEGIN J86S1
\paragraph{J86S1. A joint good eigenvector over the signature coefficient field.}
Fix an original member, a location $a$, its twist $i$, its side $b$
and its index $j=j(a)$. Fix a good row $\sigma$ with
$\mathrm{isot}_j(\sigma)\ne0$. Write $L=LL$ and $L_0=LL_0$.
The norm choice identifies the rational row space, finite action,
split operators and their eigenvalues before and after localization.
Here the group element $g$ acts in rational coordinates by
$\chi(g)^{-1}\rho(g)$. To see the identifications used below,
transport the norm choice's second-step data along its equality
with the member and its pinned character frame. Its local map reads
$\mathrm{locW}(a,\sigma,v)=R_w.\mathrm{locR}(\sigma,v)$ at split
places. Its expansion identifies the native and rational form spaces
by $e_F$ and intertwines their finite-group actions. The two
descriptions of the level subgroup agree by the member's group
reading. Both level character values have the same image under the
chosen embedding $L_0\to\operatorname{Frac}(O_{0,i})$ and its
inclusion in $\operatorname{Frac}(O_i)$, so they agree by
injectivity. In particular these are the subgroup and character in
the definition of the native space of level vectors.

Choose a nonzero $y\in\mathrm{isot}_j(\sigma)$. In the complex
target model the level projector fixes $By$ and the level subgroup
acts on it by the prescribed character. The member's extra
realization has an actual scalar system $u_{\sigma,k,v}\in L_0$:
on these complex vectors the normalized spherical double-coset sum
is multiplication by the specified embedding of $u_{\sigma,k,v}$.
The exceptional set may be taken to be the intersection of the
sets $S_0(i')$ for all twists $i'$. This is finite, being contained
in $S_0(i)$, and contains the extra realization's exceptional set
and the primes over $2g_0g_1$. For a split $v$ outside the
intersection choose $i'$ with $v\notin S_0(i')$. Conjugation
stability gives the same exclusion at its conjugate, and the pinned
geometric bad primes are contained in every $S_0(i')$. The rational
Hecke reading of the extra realization therefore applies with this
$i'$: the split diagonal with $k$ uniformizers and its right-coset
representatives realizes the stated normalized operator on rational
coordinates. The finite index of the hyperspecial intersection gives
a finite set of representatives. Its choice has no effect on
level-fixed vectors.

Consider simultaneously the kernels of the rational matrices for
the good projector minus identity, the level-group actions minus
their character, and all the split operators minus
$u_{\sigma,k,v}$. The rational target model has a basis which becomes
a basis of the complex Hilbert target; its realization in
multiplicity maps is injective. Thus $By$ gives a nonzero vector
in these kernels after extension to $\mathbb C$. In finite
dimension a finite subfamily of kernel equations cuts out their
whole intersection: take a basis of the span of all their matrix
rows. Matrix rank is unchanged by extension of fields. Consequently
some $x\ne0$ solves all the equations over $L$ in the same rational
model. The equality of the good projectors uses the pinned equality
of their group-algebra coefficients; the character twist in the
action cancels the character in the rational realization. In
particular $x$ belongs to the original rational good space, including
its level-character condition.

Put $K=\operatorname{Frac}(O_i)$ and write $\beta_b:L_0\to K$
for the side's coefficient reading. The small coefficient field
$K_s$ maps into $L_0$ by the member pin. Choose a $K_s$-basis of its
form which is also an $L$-basis after extension. Choose an
$L_0$-linear functional $f:L\to L_0$ taking a nonzero coordinate
of $x$ to $1$ (extend that coordinate to an $L_0$-basis).
Applying $\beta_b f$ to each coordinate gives a nonzero vector
in the $K$-form. Every endomorphism preserving the $K_s$-form
has matrices over $K_s$; the coordinate map therefore intertwines
it with its base change. It also transports the $L_0$-valued
character and the projector coefficients. Transport back by
$e_F^{-1}$ and call the resulting vector $z$. Then
\[
 z\ne0,\qquad e_G z=z,\qquad u z=\psi_U(u)z\quad(u\in U),\qquad
 H_{k,v} z=c_{i,b,v}^{\,k}\beta_b(u_{\sigma,k,v})z .
\]
The last identity holds for $1\le k\le r$ and every split
$v\notin S_0(i)$, with the same native unit $c_{i,b,v}$ that
relates the two Hecke actions. Each $u_{\sigma,k,v}$ used here is
the scalar in the given original realization, for both sides.
% END J86S1

% BEGIN J86S2
\paragraph{J86S2. Enumeration at this row and the Bloch--Kato good vector.}
Use $z$ from J86S1. It belongs to the level-character subspace
$F_0$, and $e_Gz=z$, so it belongs to $V_0=e_GF_0$. The
native integral submodule $B_{\rm sup}$ spans the full space over
$K$. Indeed the integral-form condition in the member's realization
places a full set of generators of the rational $K_s$-form in the
span of the integral lattice. Extending scalars gives a spanning
set for the model over $K$. The native lattice comparison contains
a common nonzero integral multiple of the transported spanning
set in $B_{\rm sup}$. Dividing by this multiplier gives the
assertion over $K$. Express $z$ as a finite linear combination
of elements of $B_{\rm sup}$; clearing the denominators gives
$0\ne d\in O_i$ with $dz\in B_{\rm sup}\cap V_0$. This intersection
is exactly the image of the native lattice $B$ in the expansion.
Thus $z$ lies in the image of $K\otimes_{O_i}B$.

On the latter space the signature span in the norm choice says that
vectors belonging to the joint eigenspaces of the signatures of this
$j,b$ span. The Hecke reading of the integral embedding intertwines
all the tested operators, also after scalar extension. Hence $z$ is
a finite sum of joint eigenvectors with systems drawn from these
native signatures. If no contributing signature shared all its
eigenvalues, for each distinct system $\lambda$ in this sum choose
one test $h$ with $\lambda(h)\ne\lambda_z(h)$. The product of
$H_h-\lambda(h)$ over these choices annihilates every summand
and is nonzero scalar multiplication on $z$. Indeed each summand
is a joint eigenvector for \emph{all} factors. This contradicts
$z\ne0$. We have found a native signature $s$ with
$j_S(i,s)=j$, $\mathrm{side}(i,s)=b$ and
\[
       h_{\rm Eig}(i,s,k,v)
             =c_{i,b,v}^{\,k}\beta_b(u_{\sigma,k,v})
             \quad\hbox{in }K .
\]
With these identifications J86S1 makes $z$ itself a witness to
the definition of $\mathrm{goodS}(i,s)$: it is nonzero, fixed by
the native good projector, and has all the specified Hecke values.
Writing $\sigma_s$ for the original row read from $s$, the norm
choice's $\mathrm{rowS\_eig}$ identifies the same left hand side
with $c_{i,b,v}^{\,k}\beta_b(u_{\sigma_s,k,v})$. Cancel the unit
and use injectivity of $\beta_b$. The two rows thus have equal
coefficients at every test for this side.
% END J86S2

% BEGIN J86S3
\paragraph{J86S3. The Galois reading of the same signature and the original triple.}
Let $\Sigma$ be the list supplied in $\mathrm{HasParamW}(a,\sigma,\Sigma)$.
The native $\mathrm{rhoS}$ specification for $s$ gives
$\mathrm{isot}_{j_S(i,s)}(\sigma_s)\ne0$ and goodness of $\sigma_s$.
Since $j_S(i,s)=j(a)$, apply U86T3--4 to $\sigma_s$ at the original
$a$. After removing the length-$t$ character string in the upper
base-change list the remaining part has length one and rank $r$.
Let $\pi$ be its base twisted by $\eta$. U86T4 gives
$\mathrm{locG}(\pi,v)=\mathrm{locW}(a,\sigma_s,v)$ at
\emph{every} split $v$. Its base is conjugate self-dual by the
base-change condition, and $\eta$ is conjugate self-dual by S86A2.
It is regular algebraic: its untwisted infinity parameters form a
submultiset of the upper list, whose exponents are distinct and lie
in $(n_V-1)/2+\mathbb Z$. The length-one part adds no Speh shifts.
The infinity exponent of $\eta$ is $1/2$ or $-1/2$ at each complex
embedding. Since $n_V$ is even and $r$ odd, shifting gives distinct
exponents in $(r-1)/2+\mathbb Z$. This proves the required test at
every embedding.

At a split place outside the finite exceptional sets the local
component of each row is spherical. The scalar for the $k$-th
hyperspecial double coset is $q_v^{k(r-k)/2}$ times the $k$-th
elementary symmetric polynomial in the Satake roots; the factor is
the square root of the modulus character of the standard Levi on
that diagonal. The measure gives the hyperspecial subgroup volume
one. The member's rational normalization multiplies each root by
$\mu_{0,j}(v)^{-1}q_v^{-n_V/2}$, and the local Frobenius reading
uses inverses of the Satake roots. Equivalently, if $A$ is the
multiset of Frobenius eigenvalues in the row reading, the
embedded rational scalar is
\[
 q_v^{k(r-k)/2}e_k\bigl(\{\alpha^{-1}\mu_{0,j}(v)^{-1}
                         q_v^{-n_V/2}:\alpha\in A\}\bigr).
\]
This follows by applying the normalized double-coset operator to a
nonzero vector of that type, using the hyperspecial fixed-vector
calculation, and cancelling the vector in the Hilbert reading.
The multisets for $\sigma,\sigma_s$ have the same size $r$: their
unramified parameters come respectively from the given rank-$r$
list and the rank-$r$ cusp $\pi$. J86S2 now equates all elementary
symmetric polynomials, since each displayed norm factor is nonzero.
Equality of monic polynomials identifies their root multisets. The
map $\alpha\mapsto\alpha^{-1}\mu_{0,j}(v)^{-1}q_v^{-n_V/2}$ is
injective on these nonzero roots. The two unramified local parameters
are therefore equal. In particular we can use the cusp $\pi$ on
the signature's row and the prescribed list on the input row even if
their finite type labels differ.

Use the reading at $a$ of $\operatorname{Frac}(O_{0,i})$ in
$\overline{\mathbb Q}_\ell$. Explicitly, extend the coefficient
map $\mathrm{twist}_2(a)$ to fractions, compose with the location's
embedding into $\overline{\mathbb Q}_\ell$, and use the corresponding
isomorphism to $\mathbb C$. The integral and side-specific
$L_0$ equalities in the localized norm choice give the identities
on $\mathbb Z_\ell$ and on $L_0$ required by $\mathrm{rhoS}$;
a fraction-field identity is checked on numerators and denominators.
Apply that specification to $\pi$. It gives semisimplicity and the
Frobenius characteristic polynomial for the twist
$\pi\otimes\eta^{-1}$ with the fixed norm shift $(1-n_V)/2$.
Semisimplicity descends through this characteristic-zero field
embedding: an invariant subspace with a complement after extension
has an equivariant retraction there. The retraction equations are
linear over the original field, and a finite spanning set of the
equations suffices, so the retraction exists over that field.

At almost every split place apply the equality of local parameters
just proved and the split-place test in $\mathrm{HasParamW}$. Twisting
each constituent of $\Sigma$ by $\eta^{-1}$ twists its full
local parameter by that character. The characteristic polynomial
for $\pi\otimes\eta^{-1}$ is therefore exactly the polynomial for
the prescribed twisted list with the same norm shift. Characteristic
polynomials commute with the coefficient embedding; on $O_{0,i}$
its composite into $\mathbb C$ is the pinned composite through
$\mathrm{twist}_2(a)$. The Frobenius assertion holds for every
valuation ring and every allowed Frobenius lift in the definition,
because the $\mathrm{rhoS}$ assertion has these quantifiers and
we have removed only a finite set of places. With the index and
side equalities from J86S2 these are all clauses of
$\mathrm{IsTriple}(a,s,\Sigma)$. The same $s$ has J86S2's
$\mathrm{goodS}$ witness. This proves the original field
$\mathrm{trip}$.
% END J86S3

\section{Multiplicity comparison for the native isotypic decomposition}

% BEGIN F86I1
\paragraph{F86I1. The parameter of a type already present at the reference level.}
Fix $\sigma$ with $n_0(\sigma)>0$. Put $G_Y=U(Y_0)$ and
$G_W=U(W_j)$ for $1\le j\le m$. We first work on $G_Y$ alone.
A reference basis row gives a nonzero vector in the level part of the
$\bar\sigma$-isotypic cuspidal multiplicity space. Complex conjugation
of both the functions and the archimedean model identifies it
antilinearly with the $\sigma$-part for the conjugate discrete series
$M_0$. Its level character is now the one read on the lift side:
the level values have modulus one. Conjugation preserves the good
projector. Indeed G85M3 pairs each finite series with its
reciprocal; the good set is closed under this operation, so their
sum of central idempotents is fixed. We thus have an actual
cuspidal automorphic representation on $G_Y$ with finite part
$\tau_f$ representing $\sigma$ and the archimedean factor prescribed
by $M_0$. This follows by decomposing the cuspidal spectrum at the
given finite level and discrete series. In the convention using
multiplicity \emph{maps} it is the archimedean representation
intertwined by those maps. The same convention is used on $G_W$.

The two real groups can have different signatures. The reference
group has signature $(k+r_S,r_S+1)$, and the $j$-th target has
signature $(k+j-1,m-j)$. Their specified real discrete series
have the same unordered Harish--Chandra parameters at each
embedding: these are the $r$ entries $\lambda_a$ of the member's
$\mathrm{lamAt}$, arranged according to the positive and negative
slots of the respective signature. In the reference model the
preceding conjugation is essential for this equality. The distinguished
place on the lift reads the original archimedean labels. The entries are integral. Here one takes the intersection with the
eventual host range for the frame and the weights, before choosing
the member. In that range $h,\Delta\ge256N(m)^4$ and the first
node $b_0\ge100r^4$. The consecutive gaps in the upper segment of
the \emph{rank-$r$} tuple are $4\Delta$, the middle gap is $b_0$,
and the later gaps are $b_{j+1}-b_j=2h$ by the node and endpoint
equations. Since $r\le N(m)$ these gaps exceed one, uniformly
in the original row index. This uses the maximum of the original
size and first-node thresholds, as in the restriction of family thresholds, and gives
a strongly regular infinitesimal character on both groups.

Use Labesse, \emph{Changement de base CM et s\'eries discr\`etes},
Corollary~5.3, on this reference representation. The CM extension is
unramified at its finite nonsplit places and $[F:\mathbb Q]>1$ as
in A72F6 (indeed its degree is even by the quadratic class
character evaluated at $-1$). Labesse's sign hypothesis $(*)$
also holds for this possibly noncompact real group. For detail,
the strongly regular case of Salamanca--Riba's unitary dual theorem
(Duke Math.\ J.\ 96 (1999), 521--546) expresses a unitary module
at this infinitesimal character as cohomological induction from a
unitary character of a $\theta$-stable Levi. The Levi blocks of
$U(p,q)$ are unitary groups. In a block of size $d$ the inducing
character is constant on its $d$ coordinates and the block's
half-sum of positive roots contributes a string with step one;
the contribution from the nilradical is constant on the block.
A block of size at least two would give two entries differing by
one in the infinitesimal character. Consequently the inducing
Levi is a compact Cartan and the resulting modules are discrete
series in its regular real packet. This is the sufficient criterion
for $(*)$ given immediately before Labesse's corollary.

The corollary gives an isobaric sum of conjugate-self-dual
discrete representations of general linear groups, with local
base-change equality at \emph{all split places} and at almost
every nonsplit place. Write the discrete factors as Speh
representations of cuspidal bases. At a designated type place
the inertia polynomial on wild invariants has degree $r$ and is
squarefree, by G85M1--3. Complex conjugation of the original
type, if used to read the polynomial, reciprocates it and
preserves both tests. The degree of a factor on wild invariants
is bounded by its rank, and the sum of all ranks with lengths
is $r$. Thus every factor has full degree there. Norm shifts
are trivial on inertia, so a length $d>1$ would repeat every
root of a nonconstant factor $d$ times. All lengths are one.
If there were more than one base, a base of smallest rank
$1\le e\le r/2$ would give a degree-$e$ divisor at every designated place.
The divisor is defined over the residue field: Frobenius permutes
the tame eigenvalues by the residue-field power, and it preserves
each Weil summand. The good row has a place admitting no such
divisor for each $1\le e\le r/2$. This proves that the base change is
one rank-$r$ cuspidal representation $\Pi$. The cuspidal clause
at the end of Labesse's corollary then gives the actual
archimedean base change of the reference discrete series, so the
real packet used below is the one for $\Pi$. The argument started
with $n_0(\sigma)>0$ and makes no assumption about the occurrence
of $\sigma$ on $W_j$.
% END F86I1

% BEGIN F86I2
\paragraph{F86I2. Multiplicity one on both inner forms and the original field \textnormal{\texttt{adm8}}.}
Use the generic pure-inner-form case of Kaletha--M\'inguez--Shin--White,
\emph{Endoscopic classification of representations: inner forms
of unitary groups}, arXiv:1409.3731, Corollary~3.3.2 and
Theorem~1.7.1 (the unconditional case specified directly after the
theorem). Both algebraic groups arise from Hermitian spaces:
an isometry over $\overline F$ to the quasi-split space gives a
Galois cocycle with values in its unitary group. Thus these are
pure inner twists. The spectral support theorem assigns an Arthur
parameter to the near-equivalence class of the reference
representation. Take the ordinary stable base-change normalization
for odd rank. The general-linear class of this parameter is determined
by its unramified transfers at almost all finite places. Labesse's
unramified comparisons in F86I1 give precisely the same class as
$\Pi$; this includes the inert as well as the split places. Uniqueness
of the expanded isobaric list from almost-everywhere agreement
(A72F5 or S86A1) identifies the parameter as the single cuspidal
$\Pi$ of Arthur length one. In particular it is generic in the
sense required in the cited theorem. Its reduced global component
group is trivial: the centralizer of a single multiplicity-one
self-dual summand consists of the two scalar signs, and both
belong to the Galois-fixed center of the dual group.
Every tensor of local packet members on either pure inner form
therefore has discrete multiplicity one.

Take on $G_W$ the finite factors of $\tau_f$ transported by
the member's finite-adelic isometry. Its conjugation of the local
groups preserves the local packet; changing the frame by an inner
automorphism gives an isomorphic representation. At infinity
choose the stipulated discrete-series model for $W_j$. The real
discrete-series correspondence for unitary groups parametrizes the
packet across signatures by the allocations of a fixed ordered
regular tuple to the positive and negative slots. F86I1 verifies
that this very model has the reference tuple, at every real
place. This proves occurrence with multiplicity one. The resulting
discrete representation is cuspidal. Indeed otherwise its
cuspidal support would be on a proper rational Levi. A proper
parabolic of a unitary group stabilizes an isotropic flag; the
dimension $e$ of its first nonzero proper step is fixed over $F$.
At every designated split place this creates a general-linear
Levi factor of rank $e$, and the local constituent is a subquotient
of induction from that Levi. The supercuspidal support of a
subquotient of such induction is the union of the induced
supports. Unramified changes in these supports do not change
their tame eigenvalues. Full wild-invariant degree therefore gives
a residue-field polynomial divisor of degree $e$ at every such
place, the same contradiction as in F86I1. This argument also
shows why residual discrete representations cannot add a copy
to the cuspidal spaces in this comparison.

It remains to compare the dimension cut out in one finite copy
on each side. The finite isometry can be chosen to match the two
lattice frames at the level primes: their integral Gram frames
agree there and the comparison has identity reduction at the
specified type primes. It carries the base subgroup and its
character to the lift subgroup and its character. On finite
vectors the other cut is the same finite sum of right translations
by lifts of the residue group, with the same good central
coefficients. Replacing a lift by another changes it by the
principal congruence kernel, which acts trivially here. The
frame used to name $\tau_f$ differs from this isometry by inner
conjugation. Thus the two cut spaces within the fixed finite
representation are isomorphic, with any conjugation of the
original reference reading already included in F86I1. The norm
factor $|P_0|$ from the realization compares squared Hilbert
norms and leaves their dimensions unchanged.

The reference $\sigma$-cut consequently has the dimension of
one copy of this finite cut. Explicitly the vectors with that
label from the orthonormal reference basis give $n_0(\sigma)$
independent vectors in it. The cut in $M_j$ contains as many
vectors, since the transported copy with the prescribed infinite
type has just been exhibited. If $n_0(\sigma)=0$ the required
lower bound is automatic. We have proved the bound in L86M5 for
every $\sigma$. Its dimension sum and independence now give
$\displaystyle M_j=\bigoplus_\sigma\mathrm{isot}_j(\sigma)$,
and in fact $\dim\mathrm{isot}_j(\sigma)=n_0(\sigma)$ for every
$\sigma$. This is the original $\mathrm{adm8}$ for all
$1\le j\le m$. Index $0$ is the reference row providing $n_0$.
% END F86I2

\section{The good lower space and multiplicities at a good point}

% BEGIN V86E1
\paragraph{V86E1. Residual congruence selects the original localization.}
At location $i$ let $\sigma$ index a nonzero summand of
$\mathrm{goodPartM}(i)$. Write $\Sigma_0,x_0$ for witnesses of
its defining condition $\mathrm{AtM}$. Apply J86O1 to the upper
theta type $P$, giving an integral character $\psi$ of the common
Hecke algebra $T_c$ for $(j(i),w(i))$. In this case $\psi$ has a
point already at $i$. To see this, take the row list $\Sigma_1$
and augmented base change of $P$ from U86T4. Both $\Sigma_0$
and $\Sigma_1$ localize to the same lower parameter at each split
place. Adding the common string and twisting therefore identifies
the local parameter of $P$ with the augmentation of $\Sigma_0$.
At a tested split place the integral polynomial $p_w$ in
$\mathrm{AtM}$ has this normalized local polynomial as its
complex image. By the coefficient calculation in J86O3, in the
common ring $O_c$ its coefficient of $X^{n_V-b}$ is
\[
    (-1)^b(\mathrm Nw)^{b(b-1)/2}\,\psi(T_{w,b})
      \qquad(1\le b\le n_V).
\]
Here we identify $O_i$ with $O_c$ by the member's
$\mathrm{phO}$, which respects the embeddings into $\mathbb C$.
The reduction of $p_w$ is that of the Hecke polynomial of $x_0$.
Since $w\nmid\ell_i$, its norm is a unit of $O_c$. Cancelling
the displayed unit for each $b$ gives equality of the two residual
characters on the generators. They also agree on constants and
hence on $T_c$. The residual kernel of $\psi$ is precisely
the maximal ideal indexing $i$, by the point description
$\mathrm{Sp.ptX}$ for $x_0$. Thus $\mathrm{Sp.ptX}$ at $i$
itself supplies $y\in X_i$ for $\psi$, since its simultaneous
eigenspace is nonzero by J86O1. The test on all Hecke generators
gives $\mathrm{ptV}(i,P)=y$.

The augmented list of $P$ contains the distinguished string,
so $y\in Z$. The test for $\mathrm{goodSig}(\Sigma_1)$ holds
because $\mathrm{HasParamW}$ identifies the row labels at the
designated split places with the list labels. Thus $y\in Z_{\rm good}$
and the original summand is in $M_{j,\Sigma_1}$ for a list above
that very point. In the opposite direction, fix $x\in Z_{\rm good}$,
any list $\Sigma$ over $x$, and any row with
$\mathrm{HasParamW}(i,\sigma,\Sigma)$. Its finite type is good
by G85M3. The list and $x$ witness $\mathrm{AtM}$: use the
Hecke polynomial of $x$ itself at each test place, since its
coefficient image equals the augmented local polynomial by the
base change of a witness of $\mathrm{HasParamX}$ and
$\mathrm{pt\_loc}$. Taking the sums of these inclusions in
both directions proves $\mathrm{RowsC.goodPart}$.
% END V86E1

% BEGIN V86E2
\paragraph{V86E2. A pointwise multiplicity test for both downward conditions.}
Fix a location $i$, a point $x\in Z_{\rm good}(i)$ and
a list $\Sigma$ over $x$. Set $A=M_{j,\Sigma}$,
$E=\mathrm{eigC}(i,x)$, and write $T$ for the original upward
theta operator. Its adjoint $T^*$ is the downward operator, with
the member's specified metrics. J86O3 gives $TA\subseteq E$,
and L86M2 gives injectivity of $T$. Also $T^*TA\subseteq A$:
on every generating lower type the composite acts by its positive
scalar by L86M4. These finite-dimensional spaces consequently satisfy
\[
 \dim A\le\dim E,\qquad E=TA\oplus\big(E\cap(TA)^\perp\big).
\]
If $\dim E\le\dim A$, equality gives $E=TA$. Then $T^*E\subseteq A$,
and for $z=Ta\in E$ the identity $T^*z=0$ implies
$\langle z,z\rangle=\langle Ta,z\rangle=0$.
Conversely, $T^*E\subseteq A$ and injectivity of $T^*$ on $E$
give $\dim E\le\dim A$. Thus the joint assertions of
$\mathrm{downMem}$ and $\mathrm{typeNZ}$ are equivalent to the
single bound
\[
                  \mathrm{mult}(i,x)\ \le\ \dim M_{j,\Sigma}
\]
for all good $x$ and every list over them. Here the identity
$\mathrm{mult}(i,x)=\dim E$ is part of the original point data.

The part of $T^*E$ detected by the named finite types can be
determined before proving this bound. If a nonzero lower type
$\sigma$ pairs with $E$ after applying $T$, its unique upper
type $P$ from U86T3 lies above $x$, by orthogonality of the upper
types. H86P3 gives equal parameters at every split place for
$P$ and a witness for $\mathrm{HasParamX}(i,x,\Sigma)$.
Untwisting and cancelling the common string in the formula of
U86T1 gives the lower localization $\mathrm{isoLoc}(\Sigma)$.
The positivity and rank clauses of $\mathrm{HasParamW}$ follow
from the augmented base-change list on removing the positive
rank-$t$ string. Hence $\sigma$ is among the summands of $A$.
If $B$ is the span of the other named lower types, this proves
$T^*E\perp B$. Moreover $TA\perp TB$ because the types lifting
from $A$ have point $x$ by J86O3 and those lifting from $B$ do
not. The positive scalar identities on the summands of $A$
give $T^*TA=A$, so also $A\perp B$. Therefore
\[
                T^*E\ \subseteq\ A\oplus (A+B)^\perp .
\]
In particular the lower spanning assertion at this row in L86M5
would give $\mathrm{downMem}$ there. For $\mathrm{typeNZ}$, the
positive scalar of $T^*T$ on the lower side controls $TA$.
The possible additional upper multiplicity is exactly
$\dim E-\dim A$. The native joint downward test requires that
this number vanish at every good point.
% END V86E2

\section{The rank-deficient averages at the designated primes}

% BEGIN N86V1
\paragraph{N86V1. Good finite series kill both mirabolic character averages.}
Fix a designated residue field $k$, a nonzero vector or covector $a$,
and the character $\theta$ of $k^\times$ prescribed by the member
there. Put $G_r=\mathrm{GL}_r(k)$. Take $H$ to be the subgroup
fixing $\ker a$ pointwise, or the subgroup acting trivially on
$k^r/ka$, respectively. For $\epsilon=1,-1$ put
\[
        m_{H,\epsilon}=\sum_{h\in H}\theta(\det h)^\epsilon h.
\]
Let $\tau$ be a finite irreducible representation in a Green
series whose polynomial omits the rank-one labels of both $\theta$
and $\theta^{-1}$. Then $m_{H,\epsilon}\tau=0$.
Indeed the normalized average is the projection onto vectors
transforming by $\theta^{-\epsilon}\circ\det$ under $H$.
In the hyperplane case, choose coordinates making $\ker a$ the first
$r-1$ axes. The subgroup consists of the unipotent radical of the
$(r-1,1)$ parabolic together with its final $\mathrm{GL}_1$ factor.
In the line case, take the line first and use the $(1,r-1)$ parabolic;
$H$ consists of its radical and its first $\mathrm{GL}_1$ factor.
In either case a nonzero projected vector lies in the radical
invariants and has the specified character on this $\mathrm{GL}_1$.
The complementary $\mathrm{GL}_{r-1}$ commutes with that factor
and acts on its character space. Finite semisimplicity then gives an
irreducible complementary constituent $\sigma$. By finite Frobenius
reciprocity, $\tau$ occurs in the Harish--Chandra induction of
$\sigma$ and $\theta^{-\epsilon}$ in the indicated order.
The cuspidal support of every such constituent contains the
rank-one label of $\theta^{-\epsilon}$, contradicting the
hypothesis. Conjugating coordinates gives every nonzero $a$.

At an original designated place the two character labels are the two
values $\chi_J,\chi_J^{-1}$ in some order, by T86C3. Each tuple
in the good idempotent $e_0$ avoids both values at each such place.
This follows also directly from the good prime-to-$\ell$ root test:
equality of a root with the forbidden value gives equality of their
prime-to-$\ell$ parts for any prime $\ell\ne\operatorname{char}k$.
Use P84C1 to expand $e_0$ as a sum of products of series idempotents.
The average in a single factor annihilates its local series idempotent
by the preceding paragraph, and commutes with all factors in other
places. Thus $m_{H,\epsilon}e_0=0$ in the actual complex group
algebra. In the member's level realization the action of the finite
quotient is by the level-character-corrected translations. Composition
of two such translations is the action of the product: the lift
discrepancy is in the kernel of reduction and its right action cancels
its level-character correction. The recorded finite-group and
coefficient identifications read the distinguished good signature as
this $e_0$ on scalar components. Hence any component fixed by it
is annihilated by all these averages. This proves the source's
$\mathrm{GoodVt}(y)\Rightarrow\mathrm{NoChi}(y)$, the
mirabolic-vanishing input in Y86N2 on the projected space. It does
not require archimedean generation or the regularized integral.
% END N86V1

\section{The pole in the downward theta integral}

% BEGIN Y86N1
\paragraph{Y86N1. The partial Euler product at the original long string.}
Let $S$ be a finite set of finite places of $F$, and omit all places
of $E$ above it in every product below. Suppose the unramified
multiset of a rank-$(r+t)$ system, after division by a unitary
character $z_w$, has a rank-$r$ remainder $C_w$ and the $t$ entries
$q_w^{(t-1)/2-j}$ for $0\le j<t$, where $q_w=\mathrm Nw$.
Assume $|b|\le q_w^{(r-1)/2}$ for $b\in C_w$ and $r+5\le t$.
This includes the source's $\mathrm{SimplePole}$ hypothesis:
there $C_w$ is the multiset $c_w$ divided by its other norm-one
scalar $m_w$. Set
\[
 L_C^S(s)=\prod_{w\nmid S}\prod_{b\in C_w}(1-bq_w^{-s})^{-1},
 \qquad s_0=(t+1)/2 .
\]
On each compact subset of $\Re s>(r+1)/2$ the logarithmic series
$\sum_{w\nmid S}\sum_{b\in C_w}\sum_{h\ge1}b^h q_w^{-hs}/h$ converges
normally. Indeed $|bq_w^{-s}|\le q_w^{-1-\delta}$ for some
$\delta>0$ there, and $\sum_w q_w^{-1-\delta}<\infty$ by the
absolutely convergent Dedekind zeta product. Its exponential is
$L_C^S(s)$, holomorphic and nowhere zero on that half-plane. For
$\Re s>s_0$ all the string factors converge absolutely as well, so
the twisted standard Euler product is
\[
 L_C^S(s)\prod_{j=0}^{t-1}\zeta_E^S(s+(t-1)/2-j).
\]
The ordering reverses the root exponents, which preserves their
multiset. The partial Dedekind zeta function is its classical
meromorphic continuation times the omitted factors
$\prod_{w\mid S}(1-q_w^{-s})$, with nonzero residue at $1$.
This yields continuation of the displayed product to
$\Re s>(r+1)/2$. At $s_0$ the zeta arguments are $t,t-1,\ldots,1$.
Exactly the last factor has a pole, since $\zeta_E^S(h)>0$ for real
$h>1$. Its product residue is therefore
\[
 L_C^S(s_0)\,\operatorname*{Res}_{s=1}\zeta_E^S(s)
        \prod_{h=2}^t\zeta_E^S(h)\ne0 .
\]
Every continuation used in the source agrees with this one by
meromorphic uniqueness from the right half-plane. This proves the
full partial-Euler-product assertion $\mathrm{SimplePole}$.

For a native good point, take a good witness list and any upper
label $P$ above the point. R86G4 gives one lower cusp; H86P3
identifies its augmented base-change list with a list for $P$
up to permutation. At every place where the list is unramified
its distinguished length-$t$ string gives the entries above after
division by the Satake value of $\mu_i\eta^{-1}$. The
remaining cusp is twisted by unitary characters only. S86A2 gives
$|b|<q_w^{1/2}\le q_w^{(r-1)/2}$, since the member has $r\ge2$.
Choose $S$ to contain the places under all ramified factors.
Both the rank guard and the pole just proved apply to every such
$S$ on this member. With the $s+1/2$ argument of the doubling
integral, the pole in its variable is $s=t/2$.
% END Y86N1

% BEGIN Y86N2
\paragraph{Y86N2. A conditional descent from the residue to the prescribed test.}
Fix a nonzero vector $y$ at an original good upper point $x$.
Use the upper realization of L86M1, and write $I_y(f,\ell,g)$ for
the downward theta integral over the member's definite source
quotient, with Schwartz function $f$, coefficient functional $\ell$
and target argument $g$. The following sufficient hypotheses
specify which assertions must hold for this same $y$.

First, suppose there are a scalar component
$u(h)=\ell(\operatorname{ext}_V(y)(h))$ and a polynomial Fock test
$\Phi$ spherical outside a finite set $S$ and principal-level
invariant at the excluded places, with the source's nonzero
archimedean integral of the Weil and automorphic matrix coefficients
for $(\Phi,u)$. Suppose $u$ is a cuspidal principal-level form with
the normalized unramified parameters of Y86N1 outside $S$.
Use the member's Weil representation, characters and Haar measures.
Suppose the doubling integral $Z_y(s)$ formed from this pair after
adjoining $t$ hyperbolic planes admits meromorphic continuation from
the region of absolute convergence to a connected region containing
$s=t/2$. Require the unramified unfolding and the excluded-place
calculation there to give, initially on a right half-plane,
\[
 Z_y(s)=C_y\,L^S(s+1/2)/d^S(s),\qquad C_y\ne0.
\]
Here the numerator is precisely the twisted product of Y86N1.
The normalizer is the product over $0\le h<a$ of
$L_F^S(2s+a-h,\eta_{E/F}^{b+h})$, with $a=r+t$, $b=r+2t$ and
$\eta_{E/F}$ the quadratic character. When $\Re s>0$ each argument
has real part greater than $1$. Absolute convergence of the logarithms
makes every factor holomorphic and nonzero near $t/2$.
Equality by continuation is taken on the component reached from the
right within the common domain and containing $t/2$.
Y86N1 then gives $\operatorname{Res}_{s=t/2} Z_y(s)\ne0$.
Assume the regularized Siegel--Weil comparison for these same data
makes this residue zero whenever all $I_y(f,\ell',g)$ vanish.
Some test, component and target argument therefore give a nonzero
integral for $y$.

For descent to the prescribed test, assume the following local
generation properties in the member's actual realization.
On the good component the inclusion of the native downward
operator $\mathrm{thG}$ equals the full fixed-test operator
$\mathrm{thF}$. Vanishing of $\mathrm{thF}(y)$ annihilates the
integrals against the designated polynomial Fock tests. Their
archimedean generation and approximation, in a topology for which
each pairing $I_y(-,\ell',g)$ is continuous, give this vanishing
on every Schwartz function whose finite Heisenberg annihilator
contains that of the member's vacuum $v_P$. Here the annihilators
consist of finitely supported group-algebra sums of Heisenberg
translations with zero archimedean component.

Require the whole Schwartz space to be spanned, allowing closure
where these pairings are continuous, by target translates of that
class, source-level character differences, and the weighted
mirabolic averages from the rank-deficient finite orbits.
The level action on $u$ for each coefficient and right translation
of Haar measure kill the differences. Require that for every
designated prime and every nonzero defining vector, the corresponding
weighted subgroup averages annihilate every scalar component of $y$.
These are the subgroups fixing the indicated hyperplane pointwise
and the subgroups acting trivially modulo the indicated line,
with both determinant-character exponents $1$ and $-1$.
Their vanishing, transported through the integral, kills the
averages. Target translations only change $g$.
Consequently $\mathrm{thG}(y)=0$ would give $I_y(f,\ell',g)=0$
for every test, every coefficient and every $g$, contradicting the
nonzero residue. If the hypotheses hold for every nonzero vector in
the indicated good eigenspace, the native downward operator is
injective there. The occurrence of a varying nonzero test establishes
only the first implication. The factorization, residue comparison
and fixed-test hypotheses above remain inputs to this criterion.
% END Y86N2

\section{The remaining multiplicity in the denominator argument}

% BEGIN F86D1
\paragraph{F86D1. One numerical test per localization.}
Fix $i$. For each good point $x$ choose a witnessing list $\Sigma_x$
and write $A_x=M_{j,\Sigma_x}$ and $E_x=\mathrm{eigC}(i,x)$.
The subspace $A_x$ is independent of the choice. In fact any two
witnesses have upper labels over the same point. Their augmented
base-change lists agree up to permutation by H86P3. Undo their common
$\eta^{-1}$ twist and subtract one copy of the marked length-$t$
term from the two multisets. The remaining lists agree up to
permutation; this preserves both the local sums and the inequalities
defining a row parameter. The same lower summands therefore generate
the two spaces.

By J86O3 the upward map $T$ takes $A_x$ into $E_x$.
The $E_x$ at different points are mutually orthogonal, since they are
sums of distinct upper types at this level, by U86T3. By V86E1 the
good row space is the sum of the $A_x$. On the entire row space $T$
is injective by L86M2. Thus
\[
 d_M(i)=\dim T(\mathrm{goodPartM}(i))
       =\sum_{x\ \mathrm{good}}\dim TA_x
       \ \le\ \sum_{x\ \mathrm{good}}\dim E_x
       =\sum_{x\ \mathrm{good}}\mathrm{mult}(i,x).
\]
All these sums are finite. Their nonnegative integer difference
is $\sum_x\dim(E_x/TA_x)$. Consequently the single reverse inequality
\[
             \sum_{x\ \mathrm{good}}\mathrm{mult}(i,x)\le d_M(i)
             \tag{2}
\]
is equivalent to $TA_x=E_x$ for every good point and every list over
it. By V86E2 this is equivalent to the joint downward membership and
downward injectivity conditions there. It also gives the original
$\mathrm{ConclusionZ}$ at $i$: the restricted kernels vanish by
injectivity of $T$, equality of the images gives the good-part
identity after taking the sum, and (2) is equality by the displayed
calculation. Conversely $\mathrm{ConclusionZ}$ supplies (2) at
nonempty good support by its last clause. On empty good support
V86E1 gives $\mathrm{goodPartM}(i)=0$, so both sides of (2) vanish
and the pointwise tests are vacuous. Thus the guard $\mathrm{HG}$
does not require an extra existence hypothesis for this equivalence.
% END F86D1

% BEGIN F86D2
\paragraph{F86D2. The remaining original member tests.}
In $\mathrm{PointConditions}'$, S86A1--4 give all ten entries of
$\mathrm{RowsAAlt}$, U86T3 gives $\mathrm{RowsB}$, and L86M2 gives
the final rank comparison for every size datum. In $\mathrm{RowsC}$,
J86O2 gives joint occurrence on the specified residual representation,
R86G7 gives the integral rank-$r$ representation, R86G4 gives the
full degree of the tame polynomial, J86O3 gives the upward containment
and V86E1 identifies the good row space. Its downward containment entry follows by the last part of V86E2
and the spanning in F86I2. Paired with the separate $\mathrm{typeNZ}$
field of $\mathrm{PointConditions}'$, it gives the two downward tests. F86D1 proves that (2) at every
$i$ is both necessary and sufficient for these two entries together.
Thus the two original fields $\mathrm{what}_2$ and $\mathrm z$ are
each equivalent on this member to (2) at all its localizations.

For $\mathrm{AdmissibleDatum}$, S86A2--4, H86P1, U86T1 and I86A3
give the local, base-change, Galois and split theta conditions.
The original large coefficient duty in M84T6 supplies
$\mathrm{hypO}$. S86A1 supplies $\mathrm{adm5}$.
The residual irreducibility condition $\mathrm{adm6}$ is the very
predicate proved in R86G6, by unfolding its definition. The last
condition is $\mathrm{adm8}$: the indicated finite types must form
an internal direct sum in every allowed lower level space.
F86I1--2 prove the missing spanning comparison for every
$1\le j\le m$ using the reference-level multiplicity and the
generic packet on the two inner forms. Together with L86M5 this
gives the internal direct sum, so the entire $\mathrm{adm}$ holds.

Sixteen of the eighteen original preconditions are now proved.
U86T4 provides
$\mathrm{Conclusion1}$, R86G6 and R86G3 give $\mathrm{rhoBar}$ and
$\mathrm{g3p}$, L86M4 supplies $\mathrm{lam}$, and J86S3 supplies
$\mathrm{trip}$. The other ten are the original fields already
established through edition 0.85. The four ties have simultaneous
witnesses by M84T10, T86C6, M84T2 and T86C4 in their defining order.
Hence (2) on the selected members supplies all the member
inputs in A85D3. Its other twist, published, finiteness, summability
and uniform family hypotheses are retained. The remaining reverse inequality (2) has been used only as a hypothesis.
% END F86D2

\section{Passing from layer points to the original family class}

% BEGIN X86B1
\paragraph{X86B1. A family bound from eventual member hypotheses.}
For a datum $D$ and its specified norm choice $a$, write
$P(D,a)$ for the conjunction of the eighteen Galois preconditions
and the four ties. Work under the published hypotheses of A85D3.
Fix a family $F$ of positive scale exponent with $\mathrm{TwGood}(F)$.
Suppose there is a real $b$ such that, for every $x\ge b$, every
$h:F.x_1<x$ and every $i\in F.I_x$,
\[
               P\bigl(F.D(x,i),F.\mathrm{ex}(x,h,i).\mathrm{nch}\bigr).
\]
Then $\mathrm{LinkN}(F)$ holds.

Put $s=\max(F.x_1,b)$ and form a family $F'$ with starting
threshold $s$. It has the original index sets at \emph{every} $x$,
including $0$, the original unbounded set of windows, the same $D,L$
and the same numerical constants. To specify every field involving
the new guard $s<x$, supply to the old field its composite
$F.x_1\le s<x$. In particular use these restrictions for the parameter
choice and its comparisons, $\mathrm{parT}$, the size and covering
properties, and $\mathrm{ex}$. All unguarded fields are the given
ones and $1<F.x_1\le s$. Equal proofs of a real inequality give
identical values of every such dependent field. This constructs
the entire $\mathrm{FamP}$ record, including its family coefficient
lists. The same witnesses and cutoff in $\mathrm{TwGood}(F)$ prove
$\mathrm{TwGood}(F')$: its norm choice and its theta datum are
the restrictions of precisely the ones used in the original test.
At every $s<x$, the displayed hypothesis now supplies the member
input of A85D3. Apply it to $F'$, with its type-prime constant $40$.

Let $C,x_0'$ be its output. Take $x_0=\max(x_0',s+1)$ on $F$.
For all $x\ge x_0$ and all $i$, both statements use the same
$D,L$, lifts, $dZ_0$, $\eta$-sum and cyclotomic series.
The $\mathrm{ShaFinOdd}$ and summability implications therefore
give exactly the desired estimate with $C,x_0$. No class membership
of $F'$ is needed in this application of the family theorem.

Consequently, for an original fixed class $\mathcal C$, it suffices
to have this eventual member assertion as a $\mathrm{Closed}$
statement on $\mathrm{Des}(\mathcal C,3)$. Retain its
$\varepsilon_1$ and its transport parameter $K(\varepsilon)$
before taking the family or $C_H$. The definition of the design
supplies $\mathrm{TwGood}$ on the given family. Choose $b$ there,
then apply the preceding argument. Separate closed inputs can be
combined with the minimum of their positive exponent cutoffs,
the pointwise supremum of the four components of their transport
parameters, and the maximum of their member thresholds on this
family. These choices retain the order of all the original quantifiers.
% END X86B1

% BEGIN X86B2
\paragraph{X86B2. Which selected row supplies the norm choice.}
Consider the predicate $P$ of X86B1. Suppose that on the original
selection $\sigma$ of a family class $\mathcal A$ we have the
following row assertion: for every $F\in\mathcal A$, every
$h:F.x_1<x$ and every $i$, the actual object
$O=\mathrm{objAt}(F,x,h,i)$ has $P(F.D(x,i),O.\mathrm{ex}.\mathrm{nch})$.
Then
\[
       \mathrm{Closed}\bigl(D_3(3,\mathcal A),\mathrm{LinkN}\bigr),
       \qquad D_3(3,\mathcal A)=
                  \mathrm{Des}(\mathrm{WithChN}(\mathcal A),3).
\]
Indeed membership projects to $F\in\mathcal A$ and to its twist
condition. The row uses the same $\mathrm{ex}$ and gives X86B1 with
$b=F.x_1$. Thus one may take $\varepsilon_1=1$ and the transport
parameter with all four components zero, uniformly before $F,C_H$.
The family output constants still have their permitted dependence.
In particular this deduction uses no separate bound for primes
$2,3$ on the selected class. The source's combination with a closed
$2,3$ row takes its outer $\varepsilon_1,K$ and applies the
family theorem using the member row alone. A85D3 already includes
the necessary arithmetic estimates on the qualifying family.

Here is a sufficient way to verify the row premise while keeping
the existing selection fixed. For each of its members require a
witness $q$ over the layer and the equality in the dependent sum
\[
        (F.D(x,i),O)=(q.w.D(q.Z),q.O).
\]
Require that its base and frame are within the range of the
layer-point member assertion, with the specified norm core, default
extension, row extension and column pin there. Evaluate the assertion
at this witness and substitute the single displayed equality:
it transports its norm argument to $O.\mathrm{ex}.\mathrm{nch}$.
By X86B1 the witnesses may also be provided for every $i$ after a
common threshold depending on the family.

The distinction between this hypothesis and the bare assertion
at layer points is visible in the source. The latter is an eventual
statement in the base context. It produces a permitted additional
condition on the slot, called \texttt{membershipLocation}, whose
content is its universally quantified conclusion at that base.
On the conjunction of this condition with the good built slot the
construction takes a built witness and its goodness, then the
division lattice and height of the selected point. Transport of the
witness along that division change produces $q$ with the displayed
equality. Its norm extension projects to the default extension,
and the witness itself gives the row extension and column pin.
This is the pointwise substitution used by the source's linked
row theorem. The linked conclusion quantifies existentially over
such an extra slot condition; the row in turn specifies an eventual
predicate $E$ and applies to selections realizing $E$. For a
previously fixed $\mathrm{Des}(\mathcal C,3)$, the specialization
must therefore supply the matching witnesses on its original
families, or show that its existing slot and selector realize
these predicates there. Inclusion of the original class in the
verified domain suffices; changing the original selected class to
a smaller class does not establish the estimate for its other
members. This describes the remaining family binding when the
remaining upper multiplicity comparison of F86D1--2 has been supplied.
% END X86B2

\endgroup
